\chapter{Viable subspaces, nets and the detectability lemma}
\label{ch:viable_subspaces}

The algorithm of Arad, Landau, Vazirani and Vidick~\cite{Arad2016RigorousRG}
keeps, for each part of a chain, a subspace of its states that, tensored with
the states of the rest of the chain, nearly contains the low-energy states;
such a subspace is called viable. This chapter proves that viability survives
tensor products, random sampling, spectral projection and the trimming of small
Schmidt coefficients, together with the nets and random projections used in the
sampling step and the detectability lemma, which bounds the energy of a state
after a product of local projections has been applied to
it~\cite{Anshu2016Simple}.

\section{Scalar soft truncation}

For $t>0$, an integer $m\geq1$, and $x\geq0$, the scalar soft truncation of
Arad--Landau--Vazirani--Vidick is the following finite logarithmic series
~\cite[lines 1152--1173]{Arad2016RigorousRG}.

\begin{definition}[Scalar soft truncation]
    \label{def:scalar_soft_truncation}
    \lean{Real.softTruncation}
    \leanok
    \begin{align}
        h_{m,t}(x)
            &=t\sum_{j=1}^{m}\frac{{(1-e^{-x/t})}^{j}}{j}.
        \label{eq:scalar_soft_truncation_definition}
    \end{align}
\end{definition}

\begin{theorem}[Scalar soft-truncation bounds]
    \label{thm:scalar_soft_truncation_bounds}
    \lean{Real.softTruncation_bounds}
    \leanok
    If $t>0$, $m\geq1$, and $x\geq0$, then
    \begin{align}
        0\leq h_{m,t}(x)&\leq x,
        \label{eq:scalar_soft_truncation_order}\\
        x-h_{m,t}(x)&\leq\frac{t}{m}{\left(\frac{x}{t}\right)}^{m},
        \label{eq:scalar_soft_truncation_remainder}\\
        h_{m,t}(x)&\leq t(1+\log m).
        \label{eq:scalar_soft_truncation_log_bound}
    \end{align}
    The final estimate corrects the printed $t\log m$ bound, which fails for
    every $m\geq1$, since $h_{m,t}(x)\to tH_m>t\log m$ as
    $x\to\infty$~\cite{gap:alvv16_soft_truncation_log_bound}.
\end{theorem}
\begin{proof}\leanok
    Put $y=x/t$, $a(y)=1-e^{-y}$, and
    $g_m(y)=\sum_{j=1}^{m}{a(y)}^{j}/j$. The finite geometric sum gives
    $g_m'(y)=1-{a(y)}^{m}$. Since $0\leq a(y)\leq\min\{1,y\}$ for $y\geq0$,
    the functions $y-g_m(y)$ and $g_m(y)-y+y^m/m$ are nondecreasing and
    vanish at zero. This proves~(\ref{eq:scalar_soft_truncation_order}) and
    (\ref{eq:scalar_soft_truncation_remainder}). Finally,
    ${a(y)}^{j}\leq1$, so the harmonic estimate
    $\sum_{j=1}^{m}1/j\leq1+\log m$ proves
    (\ref{eq:scalar_soft_truncation_log_bound}).
\end{proof}

\section{Gluing a left state through its boundary contraction}

This section records the boundary contractions and the gluing lemma of Landau,
Vazirani and Vidick~\cite[lines 409--447]{Landau2013PolynomialTime}. Fix a cut
of $(\mathcal A\otimes\mathcal S)\otimes\mathcal R$, where $\mathcal S$ is one
site, $\mathcal A$ the part strictly to its left and $\mathcal R$ the part to
its right, and let $W:\mathbb C^B\to\mathcal R$ be an isometry. In the source,
$W=U_v$ maps $\lvert j\rangle$ to the $j$-th right Schmidt vector of a vector
$v$ of Schmidt rank $B$.

\begin{definition}[Left state and boundary contraction]
    \label{def:boundary_contraction}
    \lean{Matrix.partialTraceFirst, Matrix.conjTranspose_mulVec_gluing}
    \leanok
    For a matrix $X$ on $(\mathcal A\otimes\mathcal S)\otimes\mathcal C$, write
    $\operatorname{tr}_{\mathcal A}X$ for its partial trace over
    $\mathcal A$, a matrix on $\mathcal S\otimes\mathcal C$. If
    $v=(1\otimes W)x$, then $x=(1\otimes W)^\dagger v$ is the left state of
    $v$, and $\operatorname{tr}_{\mathcal A}\lvert x\rangle\langle x\rvert$ is
    its boundary contraction.
\end{definition}

\begin{theorem}[Gluing]
    \label{thm:boundary_gluing}
    \lean{Matrix.partialTraceRight_gluing, Matrix.partialTraceFirst_gluing,
        Matrix.rectangularTraceNorm_partialTraceFirst_gluing_sub,
        Matrix.norm_trace_mul_le_l2_opNorm_mul_rectangularTraceNorm,
        Matrix.re_trace_gluing_le}
    \leanok
    \uses{def:boundary_contraction}
    Let $W$ be an isometry, $\sigma$ a matrix on
    $(\mathcal A\otimes\mathcal S)\otimes\mathbb C^B$,
    $\sigma'=(1\otimes W)\sigma(1\otimes W)^\dagger$ and $v=(1\otimes W)x$.
    Then
    \begin{enumerate}
        \item $\operatorname{tr}_{\mathcal R}\sigma'
                  =\operatorname{tr}_{\mathbb C^B}\sigma$;
        \item $\lVert\operatorname{tr}_{\mathcal A}\sigma'
                  -\operatorname{tr}_{\mathcal A}\lvert v\rangle\langle v\rvert\rVert_1
                  =\lVert\operatorname{tr}_{\mathcal A}\sigma
                  -\operatorname{tr}_{\mathcal A}\lvert x\rangle\langle x\rvert\rVert_1$;
        \item for $H=H_L\otimes1+1\otimes K$, with $H_L$ acting on
              $\mathcal A\otimes\mathcal S$ and $K$ on
              $\mathcal S\otimes\mathcal R$,
              \begin{align*}
                  \operatorname{Re}\operatorname{tr}(\sigma'H)
                   & \le\operatorname{Re}\operatorname{tr}(\sigma(H_L\otimes1))
                  +\operatorname{Re}\langle v,(1\otimes K)v\rangle
                  +\lVert K\rVert\,\lVert\operatorname{tr}_{\mathcal A}\sigma
                  -\operatorname{tr}_{\mathcal A}\lvert x\rangle\langle x\rvert\rVert_1.
              \end{align*}
    \end{enumerate}
    This is Lemma~l:gluing of~\cite[lines 421--447]{Landau2013PolynomialTime},
    for every isometry $W$ whose range contains the right Schmidt vectors of
    $v$. The constant $n$ of the printed third part is replaced by $\lVert K\rVert$,
    which is at most $n$ for the chain Hamiltonian of the source, where
    $K=H_i+H_R$.
\end{theorem}

\begin{proof}
    \leanok
    The first part holds because $W^\dagger W=1$. For any matrix $X$,
    $\operatorname{tr}_{\mathcal A}((1\otimes W)X(1\otimes W)^\dagger)
        =(1\otimes W)(\operatorname{tr}_{\mathcal A}X)(1\otimes W)^\dagger$, and
    $\lvert v\rangle\langle v\rvert
        =(1\otimes W)\lvert x\rangle\langle x\rvert(1\otimes W)^\dagger$; the
    second part follows since conjugation by an isometry preserves the trace
    norm. For the third, $(1\otimes W)^\dagger(H_L\otimes1)(1\otimes W)
        =H_L\otimes1$, and
    $\operatorname{tr}(\sigma'(1\otimes K))
        -\langle v,(1\otimes K)v\rangle
        =\operatorname{tr}\bigl((\operatorname{tr}_{\mathcal A}\sigma'
        -\operatorname{tr}_{\mathcal A}\lvert v\rangle\langle v\rvert)K\bigr)$,
    whose modulus is at most $\lVert K\rVert$ times the trace norm of the
    difference, which is the quantity of the second part.
\end{proof}

\begin{theorem}[A dominant low-energy eigenvector]
    \label{thm:dominant_low_energy_eigenvector}
    \lean{Matrix.IsHermitian.re_trace_mul_eq_sum_eigenvalues_mul,
        Matrix.PosSemidef.exists_eigenvector_of_re_trace_mul_le}
    \leanok
    Let $H-E_0\ge\Delta(1-\lvert\Gamma\rangle\langle\Gamma\rvert)$ with
    $\Delta>0$ and $\lVert\Gamma\rVert=1$, and let $\sigma$ be a density matrix
    with $\operatorname{Re}\operatorname{tr}(\sigma H)\le E_0+e$, where
    $4e<\Delta$. Then $\sigma$ has a unit eigenvector $v$ with eigenvalue larger
    than $\tfrac12$, energy
    $\operatorname{Re}\langle v,Hv\rangle\le E_0+2e$ and overlap
    $\lvert\langle\Gamma,v\rangle\rvert\ge1-2e/\Delta$.
\end{theorem}

\begin{proof}
    \leanok
    Write $\sigma=\sum_k\lambda_k\lvert v_k\rangle\langle v_k\rvert$, so that
    $\sum_k\lambda_k(\langle v_k,Hv_k\rangle-E_0)\le e$ with every term
    nonnegative. The eigenvectors of energy above $E_0+2e$ therefore carry total
    weight less than $\tfrac12$. Each eigenvector of energy at most $E_0+2e$ has
    $\lvert\langle\Gamma,v_k\rangle\rvert^2\ge1-2e/\Delta>\tfrac12$, and since
    $\sum_k\lvert\langle\Gamma,v_k\rangle\rvert^2=1$ there is at most one of
    them; its eigenvalue exceeds $\tfrac12$, and the overlap is at least its
    square.
\end{proof}

\begin{theorem}[The leading eigenvector of a near-optimal left state]
    \label{thm:leading_eigenvector_gluing}
    \lean{Matrix.trace_gluing_mul_kronecker_one,
        Matrix.exists_leading_eigenvector_of_gluing}
    \leanok
    \uses{thm:boundary_gluing, thm:dominant_low_energy_eigenvector}
    In the setting of Theorem~\ref{thm:boundary_gluing}, let
    $H-E_0\ge\Delta(1-\lvert\Gamma\rangle\langle\Gamma\rvert)$ with
    $\Delta>0$ and $\lVert\Gamma\rVert=1$, and let $v=(1\otimes W)x$ have
    $\operatorname{Re}\langle v,Hv\rangle\le E_0+e$. Let $\sigma$ be a density
    matrix on $(\mathcal A\otimes\mathcal S)\otimes\mathbb C^B$ with
    $\operatorname{Re}\operatorname{tr}(\sigma(H_L\otimes1))\le
        \operatorname{Re}\operatorname{tr}(\lvert x\rangle\langle x\rvert(H_L\otimes1))$
    and
    $\lVert\operatorname{tr}_{\mathcal A}\sigma
        -\operatorname{tr}_{\mathcal A}\lvert x\rangle\langle x\rvert\rVert_1\le\tau$.
    If $4(e+\lVert K\rVert\tau)<\Delta$, then $\sigma$ has a unit eigenvector
    $u$ with eigenvalue larger than $\tfrac12$ such that $u'=(1\otimes W)u$ has
    energy at most $E_0+2(e+\lVert K\rVert\tau)$ and
    $\lvert\langle\Gamma,u'\rangle\rvert\ge1-2(e+\lVert K\rVert\tau)/\Delta$.

    This is Lemma~l:leading of~\cite[lines 493--505]{Landau2013PolynomialTime},
    for the solution $\sigma$ of the size-trimming program: its left state $x$ of
    $v$ is feasible, so $\sigma$ has left energy at most that of $x$, and
    $\tau=c_\varepsilon/n$. With $\lVert K\rVert\le n$, $e$ is the energy
    excess of the trimmed state of Lemma~l:within, whose printed value
    $6\sqrt{c_\varepsilon}$ is not established by the printed argument; the
    corrected excess and a choice of $c_\varepsilon$ with
    $2(e+c_\varepsilon)<\varepsilon/12$ are given
    in~\cite{gap:lvv13_trimming_energy}.
\end{theorem}

\begin{proof}
    \leanok
    By Theorem~\ref{thm:boundary_gluing} and the left-energy bound,
    $\sigma'=(1\otimes W)\sigma(1\otimes W)^\dagger$ has energy at most
    $E_0+e+\lVert K\rVert\tau$, and it is a density matrix. Theorem~\ref{thm:dominant_low_energy_eigenvector}
    gives an eigenvector $v'$ of $\sigma'$ with eigenvalue $\mu>\tfrac12$. Since
    $\mu\ne0$, $v'=\mu^{-1}\sigma'v'$ lies in the range of $1\otimes W$, so
    $u=(1\otimes W)^\dagger v'$ satisfies $(1\otimes W)u=v'$ and
    $\sigma u=\mu u$.
\end{proof}

\section{A finite net for the trace-norm unit ball}

The cardinality reduction of Landau, Vazirani and
Vidick~\cite[lines 451--460]{Landau2013PolynomialTime} discretizes boundary
contractions by a finite net in trace norm.

\begin{theorem}[Entries and trace norm]
    \label{thm:trace_norm_entries}
    \lean{Matrix.norm_apply_le_l2_opNorm, Matrix.norm_apply_le_rectangularTraceNorm,
        Matrix.rectangularTraceNorm_single_one_le,
        Matrix.rectangularTraceNorm_le_sum_norm_apply}
    \leanok
    For a complex square matrix $Y$, every entry satisfies
    $\lvert Y_{ij}\rvert\le\lVert Y\rVert\le\lVert Y\rVert_1$, and
    $\lVert Y\rVert_1\le\sum_{i,j}\lvert Y_{ij}\rvert$.
\end{theorem}

\begin{proof}
    \leanok
    An entry is bounded by the norm of its column, which is at most the
    operator norm. The entry $Y_{ij}$ is the trace pairing of $Y$ with a
    matrix unit, which has operator norm one, so it is at most
    $\lVert Y\rVert_1$. Writing $Y$ as the sum of its entries times matrix
    units, each of trace norm at most one, gives the last bound.
\end{proof}

\begin{theorem}[A trace-norm net]
    \label{thm:trace_norm_net}
    \lean{Matrix.exists_finset_rectangularTraceNorm_net}
    \leanok
    \uses{thm:trace_norm_entries}
    Let $c$ be the size of a square matrix and $\eta>0$. There is a set of at
    most $(2\lceil c^2/\eta\rceil+1)^{2c^2}$ complex $c\times c$ matrices such
    that every matrix of trace norm at most $1$ is within trace-norm distance
    $\eta$ of one of them. This is Lemma~l:constantbondapprox
    of~\cite[lines 451--460]{Landau2013PolynomialTime} for $c=Bd$, with the
    printed grid step, net and size corrected as recorded
    in~\cite{gap:lvv13_trace_norm_net}.
\end{theorem}

\begin{proof}
    \leanok
    Put $M=\lceil c^2/\eta\rceil$ and round the real and imaginary part of
    every entry, which lies in $[-1,1]$, to the grid
    $\{-1+k/M:0\le k\le2M\}$. Each entry moves by at most $1/M$, so the trace
    norm of the difference is at most $c^2/M\le\eta$, and there are at most
    $(2M+1)^{2c^2}$ rounded matrices.
\end{proof}

\section{Nets of the unit ball by volume}

Let $E$ be a real normed space of dimension $n$.

\begin{theorem}[Volume bound on separated sets]
    \label{thm:ball_separated_card}
    \lean{Metric.card_le_of_le_dist_of_subset_closedBall}
    \leanok
    A finite subset of the closed unit ball of $E$ whose distinct points are
    at distance at least $\varepsilon>0$ has at most $(1+2/\varepsilon)^n$
    elements.
\end{theorem}

\begin{proof}
    \leanok
    The open balls of radius $\varepsilon/2$ about the points are disjoint
    and lie in the ball of radius $1+\varepsilon/2$ about the origin. Haar
    measure multiplies the volume of a ball by $r^n$ when its radius is
    multiplied by $r$, so the number of points times $(\varepsilon/2)^n$ is
    at most $(1+\varepsilon/2)^n$.
\end{proof}

\begin{theorem}[Packing bound]
    \label{thm:ball_packing}
    \lean{Metric.packingNumber_le_of_subset_closedBall}
    \leanok
    \uses{thm:ball_separated_card}
    For $\varepsilon>0$, every set of points of a subset of the closed unit
    ball of $E$ at pairwise distance more than $\varepsilon$ has at most
    $(1+2/\varepsilon)^n$ elements.
\end{theorem}

\begin{proof}
    \leanok
    An infinite or larger such set would contain a finite subset with more
    elements than the bound, contradicting the previous theorem.
\end{proof}

\begin{theorem}[Volume net]
    \label{thm:ball_net}
    \lean{Metric.exists_finset_net_of_subset_closedBall}
    \leanok
    \uses{thm:ball_packing}
    Every subset $A$ of the closed unit ball of $E$ contains a finite set $C$
    of at most $(1+2/\varepsilon)^n$ points such that every point of $A$ lies
    within distance $\varepsilon$ of a point of $C$. This is the volume
    argument used in the proof of Lemma~l:rs
    of~\cite[lines 541--546]{Arad2016RigorousRG}; for a complex space of
    dimension $r$, $n=2r$.
\end{theorem}

\begin{proof}
    \leanok
    Take a maximal subset $C$ of $A$ whose points are at pairwise distance
    more than $\varepsilon$; it is finite with the stated size by the packing
    bound. A point of $A$ at distance more than $\varepsilon$ from every point
    of $C$ could be added to $C$, so there is none.
\end{proof}

\section{Projection onto a Haar-random coordinate subspace}

The random sampling lemma of Arad, Landau, Vazirani and
Vidick~\cite[lines 532--547]{Arad2016RigorousRG} bounds the squared norm of the
projection of a fixed unit vector onto a uniformly random subspace by a
standard concentration argument. The identity below reduces that random
projection to independent Gaussian coordinates.

Let $\iota$ be a finite index set, let $A\subseteq\iota$, and let $W_0$ be the
span of the standard basis vectors $e_i$ with $i\in A$. Let $U$ be distributed
by the Haar probability measure on the unitary group of $\mathbb{C}^\iota$, and
let $(g_i)_{i\in\iota}$ be independent standard circular complex Gaussians,
$g_i=(a_i+\mathrm{i}\,b_i)/\sqrt2$ with $a_i,b_i$ independent standard real
Gaussians.

\begin{theorem}[Projection onto a rotated coordinate subspace]
    \label{thm:haar_projection_coordinate_weight}
    \lean{QICLean.HaarProjection.starProjection_coordinateSubspace,
        QICLean.HaarProjection.norm_sq_starProjection_map_coordinateSubspace}
    \leanok
    For every unitary $U$ and every vector $x$,
    $\lVert P_{UW_0}x\rVert^2=\sum_{i\in A}\lvert(U^\dagger x)_i\rvert^2$.
\end{theorem}

\begin{proof}
    \leanok
    The projection onto $W_0$ keeps the coordinates in $A$ and discards the
    others, and $P_{UW_0}=U P_{W_0}U^\dagger$ with $U$ preserving norms.
\end{proof}

\begin{theorem}[Law of the projection onto a Haar-random subspace]
    \label{thm:haar_projection_law}
    \lean{QICLean.HaarProjection.unitaryHaar_sum_norm_sq_le,
        QICLean.HaarProjection.unitaryHaar_norm_sq_starProjection_le}
    \leanok
    \uses{thm:haar_projection_coordinate_weight}
    For every unit vector $x\in\mathbb{C}^\iota$ and every real $c$,
    \[
        \Pr_U\Bigl[\lVert P_{UW_0}x\rVert^2\le c\Bigr]
        =\Pr_g\Bigl[\sum_{i\in A}\lvert g_i\rvert^2
            \le c\sum_{i\in\iota}\lvert g_i\rvert^2\Bigr].
    \]
\end{theorem}

\begin{proof}
    \leanok
    Let $K$ be the set of vectors $y$ with
    $\sum_{i\in A}\lvert y_i\rvert^2\le c\sum_i\lvert y_i\rvert^2$; it is
    closed and invariant under multiplication by nonzero complex scalars. By
    the lemma, and since $\lVert U^\dagger x\rVert=1$, the left side is the
    Haar probability of $U^\dagger x\in K$. The vector $g$ is $2^{-1/2}$ times
    a standard Gaussian vector $y$ of $\mathbb{C}^\iota$ regarded as a real
    space, so the right side is the probability of $y\in K$.

    Integrate the indicator of $U^\dagger y\in K$ against the product of the
    Haar measure and the law of $y$ in both orders. For fixed $U$, the law of
    $y$ is invariant under the real isometry $U^\dagger$, so the inner integral
    is the probability of $y\in K$. For fixed $y\neq0$, choose a unitary $V$
    with $Vx=y/\lVert y\rVert$; then
    $U^\dagger y=\lVert y\rVert\,(V^\dagger U)^\dagger x$, and invariance of
    the Haar measure under left multiplication by $V^\dagger$ shows that the
    inner integral is the Haar probability of $U^\dagger x\in K$. The origin is
    a null set for the Gaussian, so both iterated integrals agree with the two
    sides of the identity.
\end{proof}

\section{Closeness and viability of subspaces}

Arad, Landau, Vazirani and Vidick~\cite[Section~s:vs]{Arad2016RigorousRG}
measure how well one subspace approximates another through orthogonal
projections. For a subspace $S$ of one tensor factor, the extended subspace
$S\otimes\mathcal H_B$ plays the role of $T$ below, and $\delta$-viability of
$S$ for $T'$ is $\delta$-closeness of $S\otimes\mathcal H_B$ to $T'$.

\begin{definition}[Closeness]
    \label{def:subspace_close}
    \lean{Submodule.IsCloseTo}
    \leanok
    A subspace $T$ is $\delta$-close to a subspace $T'$ if
    $P_{T'}P_TP_{T'}\ge(1-\delta)P_{T'}$, where $P_T$ and $P_{T'}$ are the
    orthogonal projections~\cite[lines 433--440 and 469--474]{Arad2016RigorousRG}.
\end{definition}

\begin{theorem}[Closeness in terms of vectors]
    \label{thm:subspace_close_iff}
    \lean{Submodule.isCloseTo_iff}
    \leanok
    \uses{def:subspace_close}
    $T$ is $\delta$-close to $T'$ exactly when
    $\lVert P_Tt\rVert^2\ge(1-\delta)\lVert t\rVert^2$ for every $t\in T'$.
\end{theorem}

\begin{proof}
    \leanok
    The quadratic form of $P_{T'}P_TP_{T'}-(1-\delta)P_{T'}$ at $x$ is
    $\lVert P_TP_{T'}x\rVert^2-(1-\delta)\lVert P_{T'}x\rVert^2$, and
    $P_{T'}x$ ranges over $T'$.
\end{proof}

\begin{theorem}[Preimages under the projection]
    \label{thm:subspace_close_preimage}
    \lean{Submodule.IsCloseTo.exists_mem_starProjection_eq}
    \leanok
    \uses{thm:subspace_close_iff}
    If $S$ is $\delta$-close to $T$ with $\delta<1$, then every $t\in T$ is
    $P_Ts$ for some $s\in S$ with
    $\lVert s\rVert\le\lVert t\rVert/(1-\delta)$. This is Lemma~l:2
    of~\cite[lines 449--465]{Arad2016RigorousRG}.
\end{theorem}

\begin{proof}
    \leanok
    The map $r\mapsto P_TP_Sr$ of $T$ into itself satisfies
    $\operatorname{Re}\langle P_TP_Sr,r\rangle=\lVert P_Sr\rVert^2
    \ge(1-\delta)\lVert r\rVert^2$, so it is injective, hence surjective in
    finite dimension. Take $r$ with $P_TP_Sr=t$ and $s=P_Sr$. Then
    $(1-\delta)\lVert r\rVert^2\le\operatorname{Re}\langle t,r\rangle
    \le\lVert t\rVert\,\lVert r\rVert$, and $\lVert s\rVert\le\lVert r\rVert$.
\end{proof}

\begin{theorem}[Robustness]
    \label{thm:subspace_close_trans}
    \lean{Submodule.IsCloseTo.trans}
    \leanok
    \uses{thm:subspace_close_iff}
    If $T$ is $\delta$-close to $T'$ with $\delta\ge0$ and $T'$ is
    $\delta'$-close to $T''$, then $T$ is $2(\delta+\delta')$-close to $T''$.
    Taking $T=S\otimes\mathcal H_B$, a $\delta$-viable set for $T'$ is
    $2(\delta+\delta')$-viable for $T''$. This is Lemma~l:close
    of~\cite[lines 484--489]{Arad2016RigorousRG}.
\end{theorem}

\begin{proof}
    \leanok
    For $a\in T''$ let $b=P_{T'}a$ and $c=P_Tb$. By Pythagoras
    $\lVert a-b\rVert^2\le\delta'\lVert a\rVert^2$ and
    $\lVert b-c\rVert^2\le\delta\lVert b\rVert^2\le\delta\lVert a\rVert^2$, so
    $\lVert a-c\rVert^2\le2(\delta+\delta')\lVert a\rVert^2$. Since $c\in T$,
    $\lVert a-P_Ta\rVert\le\lVert a-c\rVert$, and Pythagoras for $P_T$ gives
    the claim.
\end{proof}

\begin{theorem}[Projection onto an intersection]
    \label{thm:subspace_inf_projection}
    \lean{Submodule.starProjection_inf_of_commute}
    \leanok
    If the orthogonal projections onto $U_1$ and $U_2$ commute, the
    projection onto $U_1\cap U_2$ is $P_{U_1}P_{U_2}$.
\end{theorem}

\begin{proof}
    \leanok
    The vector $P_{U_1}P_{U_2}x=P_{U_2}P_{U_1}x$ lies in both subspaces, and
    $x-P_{U_1}P_{U_2}x$ is orthogonal to every $w\in U_1\cap U_2$, since
    $\langle P_{U_1}P_{U_2}x,w\rangle=\langle x,P_{U_2}P_{U_1}w\rangle
    =\langle x,w\rangle$.
\end{proof}

\begin{theorem}[Closeness of an intersection]
    \label{thm:subspace_close_inf}
    \lean{Submodule.IsCloseTo.inf_of_commute}
    \leanok
    \uses{thm:subspace_close_iff, thm:subspace_inf_projection}
    If $U_1$ is $\delta_1$-close and $U_2$ is $\delta_2$-close to $T$, and
    the projections onto $U_1$ and $U_2$ commute, then $U_1\cap U_2$ is
    $(\delta_1+\delta_2)$-close to $T$. This is the operator step in the
    proof of Lemma~l:t of~\cite[lines 506--523]{Arad2016RigorousRG}: for
    sets $S_1$, $S_2$ on disjoint subsystems, the extensions
    $U_k=S_k^{(\mathrm{ext})}$ have commuting projections and
    $U_1\cap U_2=(S_1\otimes S_2)^{(\mathrm{ext})}$. The tensor-product
    statement of Lemma~l:t is not formalized here.
\end{theorem}

\begin{proof}
    \leanok
    For $t\in T$ write $P_1t=P_1P_2t+P_1(t-P_2t)$. The two summands are
    orthogonal, so
    $\lVert P_1P_2t\rVert^2=\lVert P_1t\rVert^2-\lVert P_1(t-P_2t)\rVert^2
    \ge(1-\delta_1)\lVert t\rVert^2-\lVert t-P_2t\rVert^2$, and
    $\lVert t-P_2t\rVert^2=\lVert t\rVert^2-\lVert P_2t\rVert^2
    \le\delta_2\lVert t\rVert^2$.
\end{proof}

\begin{theorem}[Dimension of an approximated subspace]
    \label{thm:subspace_close_dimension}
    \lean{Submodule.IsCloseTo.finrank_le}
    \leanok
    \uses{def:subspace_close}
    Let $U$ and $T$ be subspaces of a finite-dimensional real or complex
    inner-product space. If $U$ is $\delta$-close to $T$ and $\delta<1$,
    then $\dim T\le\dim U$. This follows from the projection preimage
    estimate of~\cite[lines 449--465]{Arad2016RigorousRG}.
\end{theorem}

\begin{proof}
    \leanok
    \uses{thm:subspace_close_preimage}
    The orthogonal projection $P_T$ maps $U$ onto $T$, so the dimension
    inequality follows from surjectivity.
\end{proof}

\begin{theorem}[Equality of nested approximating subspaces]
    \label{thm:subspace_close_nested_equal}
    \lean{Submodule.eq_and_isCloseTo_of_le}
    \leanok
    \uses{def:subspace_close}
    Let $U\subseteq V$ and $G$ be subspaces of a finite-dimensional real
    or complex inner-product space. Suppose $U$ is $\delta$-close to $G$
    and $G$ is $\delta'$-close to $V$, where $\delta,\delta'<1$.
    Then $U=V$. Every subspace $W$ with $U\subseteq W\subseteq V$ equals
    $U$, is $\delta$-close to $G$, and is approximated by $G$ with error
    $\delta'$. This is the dimension comparison used
    in~\cite[lines 1010--1024]{Arad2016RigorousRG}.
\end{theorem}

\begin{proof}
    \leanok
    \uses{thm:subspace_close_dimension}
    The two approximation bounds and the inclusion give
    $\dim G\le\dim U\le\dim V\le\dim G$. Hence $U=V$, and any intermediate
    subspace is the same space and retains both approximation bounds.
\end{proof}

\section{The lower tail of a complex Gaussian coordinate ratio}

The random sampling lemma of Arad, Landau, Vazirani and
Vidick~\cite[lines 531--547]{Arad2016RigorousRG} selects a random
$s$-dimensional subspace of a $q$-dimensional space and needs that a fixed
vector keeps at least half of its expected squared weight on it, except with
probability $e^{-s/16}$. The source quotes this from the real Gaussian
estimate of Dasgupta and Gupta~\cite[Theorem~2.1]{Dasgupta2003ElementaryJL}.
For standard circular complex Gaussian coordinates the bound follows from a
Chernoff estimate with one fixed parameter, as below. Throughout,
$g_1,\dots,g_q$ are independent standard circular complex Gaussians, each the
sum $(x+\mathrm{i}y)/\sqrt2$ of two independent standard real Gaussians, so
that $\mathbb E\lvert g_i\rvert^2=1$.

\begin{theorem}[Exponential moments of the squared moduli]
    \label{thm:complex_gaussian_norm_sq_mgf}
    \lean{QICLean.ComplexGaussian.integral_exp_sum_mul_norm_coordinate_sq,
        QICLean.ComplexGaussian.integral_exp_mul_norm_coordinate_sq}
    \leanok
    For real weights $c_1,\dots,c_q<1$,
    \[
        \mathbb E\exp\Big(\sum_{i=1}^q c_i\lvert g_i\rvert^2\Big)
        =\prod_{i=1}^q\frac1{1-c_i}.
    \]
    In particular $\mathbb E\exp(c\lvert g_i\rvert^2)=1/(1-c)$ for $c<1$.
\end{theorem}

\begin{proof}
    \leanok
    Write $\lvert g_i\rvert^2=(x_i^2+y_i^2)/2$. The exponential factors into
    one factor for each of the $2q$ independent real coordinates, and for a
    standard real Gaussian $t$ and $b<1/2$ the Gaussian integral gives
    $\mathbb E e^{bt^2}=(1-2b)^{-1/2}$. With $b=c_i/2$ the two factors of the
    $i$-th coordinate multiply to $1/(1-c_i)$.
\end{proof}

\begin{theorem}[Lower tail of the coordinate ratio]
    \label{thm:complex_gaussian_ratio_lower_tail}
    \lean{QICLean.ComplexGaussian.measureReal_sum_norm_coordinate_sq_le_le,
        QICLean.ComplexGaussian.measure_sum_norm_coordinate_sq_le_le}
    \leanok
    \uses{thm:complex_gaussian_norm_sq_mgf}
    Let $A\subseteq\{1,\dots,q\}$ have $s$ elements. Then
    \[
        \Pr\Big(\sum_{i\in A}\lvert g_i\rvert^2
        \le\frac{s}{2q}\sum_{i=1}^q\lvert g_i\rvert^2\Big)\le e^{-s/16}.
    \]
\end{theorem}

\begin{proof}
    \leanok
    Put $a=s/(2q)\le1/2$ and $w_i=1-a$ for $i\in A$, $w_i=-a$ otherwise. The
    event is $Z\le0$ for $Z=\sum_iw_i\lvert g_i\rvert^2$, so Markov's
    inequality for $e^{-Z/2}$ and the exponential moments above bound its
    probability by $\prod_i(1+w_i/2)^{-1}$. Each factor $u$ satisfies
    $u\le e^{u-1}$, so the probability is at most $e^E$ with
    \[
        E=-s\,\frac{2q-s}{6q-s}+(q-s)\frac{s}{4q-s}.
    \]
    For $0<s\le q$, the inequality $E\le-s/16$ is equivalent, after
    multiplying by $16(6q-s)(4q-s)/s>0$, to $s^2-26qs-8q^2\le0$, which holds
    since $s\le q$. For $s=0$ both sides of the claim are trivial.
\end{proof}

\section{Error reduction of close subspaces}

Arad, Landau, Vazirani and Vidick~\cite[Lemma~l:er]{Arad2016RigorousRG}
improve a viable set by applying an approximate ground state projection.
The operator step of their argument holds for any operator that preserves
the target subspace and its orthogonal complement, does not shrink vectors
of the target, and contracts its complement.

\begin{theorem}[Monotonicity of closeness]
    \label{thm:subspace_close_mono}
    \lean{Submodule.IsCloseTo.mono}
    \leanok
    \uses{thm:subspace_close_iff}
    If $U$ is $\delta$-close to $T$ and $U\subseteq U'$, then $U'$ is
    $\delta$-close to $T$.
\end{theorem}

\begin{proof}
    \leanok
    Since $P_U=P_UP_{U'}$, we have
    $\lVert P_Ut\rVert\le\lVert P_{U'}t\rVert$ for every $t$.
\end{proof}

\begin{theorem}[Error reduction by a contraction]
    \label{thm:subspace_close_contract}
    \lean{Submodule.IsCloseTo.map_of_contract}
    \leanok
    \uses{thm:subspace_close_iff, thm:subspace_close_preimage}
    Let $K$ map $T$ and $T^\perp$ into themselves, with
    $\lVert Kt\rVert\ge\lVert t\rVert$ for $t\in T$ and
    $\lVert Kq\rVert^2\le\Delta\lVert q\rVert^2$ for $q\in T^\perp$, where
    $\Delta\ge0$. If $U$ is $\delta$-close to $T$ with $\delta<1$, then $KU$
    is $\Delta/(1-\delta)^2$-close to $T$. This is the operator step in the
    proof of Lemma~l:er of~\cite[lines 606--622]{Arad2016RigorousRG}: for
    $K=\sum_iL_i\otimes A_i\otimes R_i$ and
    $U=\mathcal H_L\otimes S\otimes\mathcal H_R$, the image $KU$ lies in
    $\mathcal H_L\otimes V\otimes\mathcal H_R$ with
    $V=\operatorname{span}\{A_iS\}$. The tensor-product statement of
    Lemma~l:er, with $T$ a spectral subspace of a Hamiltonian, is not
    formalized here.
\end{theorem}

\begin{proof}
    \leanok
    For $t\in T$, the restriction of $K$ to $T$ is injective, hence
    surjective, so $t=Kr$ with $r\in T$ and $\lVert r\rVert\le\lVert t\rVert$.
    Take $u\in U$ with $P_Tu=r$ and
    $\lVert u\rVert\le\lVert r\rVert/(1-\delta)$, and put $q=u-r\in T^\perp$.
    Then $Ku=t+Kq$ with orthogonal summands, and
    $y=\lVert Kq\rVert^2\le\Delta\lVert u\rVert^2
    \le\Delta\lVert t\rVert^2/(1-\delta)^2$. Since $Ku\in KU$,
    $\lVert t\rVert^2=\langle Ku,P_{KU}t\rangle
    \le\lVert Ku\rVert\,\lVert P_{KU}t\rVert$, so with
    $x=\lVert t\rVert^2$ we get
    $\lVert P_{KU}t\rVert^2\ge x^2/(x+y)\ge x-y$.
\end{proof}

\section{Viable subspaces of a tensor product}

For subspaces $S\subseteq E$ and $S'\subseteq F$ of finite-dimensional
inner product spaces, $S\otimes S'$ denotes the image of $S\otimes S'$ in
$E\otimes F$ under the tensor product of the two inclusions.

\begin{theorem}[Projection onto a tensor product of subspaces]
    \label{thm:subspace_tensor_projection}
    \lean{Submodule.starProjection_tensorProd}
    \leanok
    The orthogonal projection onto $S\otimes S'$ is $P_S\otimes P_{S'}$.
\end{theorem}

\begin{proof}
    \leanok
    The operator $P_S\otimes P_{S'}$ is idempotent and self-adjoint, so it
    is the orthogonal projection onto its range. Writing $P_S$ as the
    inclusion of $S$ composed with the surjection $E\to S$, and likewise
    for $S'$, the range is the image of $S\otimes S'$.
\end{proof}

\begin{definition}[Viable subspace]
    \label{def:subspace_viable}
    \lean{Submodule.IsViable}
    \leanok
    \uses{def:subspace_close}
    A subspace $S\subseteq E$ is $\delta$-viable for $T\subseteq E\otimes F$
    if $S\otimes F$ is $\delta$-close to $T$, that is,
    $P_TP_{S\otimes F}P_T\ge(1-\delta)P_T$. This is Definition~d:vs
    of~\cite[lines 433--440]{Arad2016RigorousRG}.
\end{definition}

\begin{theorem}[Tensoring]
    \label{thm:viable_tensoring}
    \lean{Submodule.IsViable.tensorProd}
    \leanok
    \uses{def:subspace_viable, thm:subspace_tensor_projection, thm:subspace_inf_projection, thm:subspace_close_inf}
    Let $T\subseteq E\otimes F\otimes G$. If $S_1\otimes F$ is
    $\delta_1$-viable and $E\otimes S_2$ is $\delta_2$-viable for $T$, then
    $S_1\otimes S_2$ is $(\delta_1+\delta_2)$-viable for $T$. Equivalently,
    viable sets $S_1\subseteq E$ and $S_2\subseteq F$ on disjoint
    subsystems tensor to a $(\delta_1+\delta_2)$-viable set. This is
    Lemma~l:t of~\cite[lines 500--523]{Arad2016RigorousRG}.
\end{theorem}

\begin{proof}
    \leanok
    The projections onto $S_1\otimes F\otimes G$ and $E\otimes S_2\otimes G$
    are $P_{S_1}\otimes1\otimes1$ and $1\otimes P_{S_2}\otimes1$. They
    commute, and their product $P_{S_1}\otimes P_{S_2}\otimes1$ is the
    projection onto $S_1\otimes S_2\otimes G$, which is therefore the
    intersection of the two subspaces. The closeness of an intersection
    gives the claim.
\end{proof}

\section{Viable subspaces in coordinates}

When $\mathcal H_A=\mathbb C^\sigma$ is the space of configurations
$\sigma$ of a finite set of sites and $\mathcal H_B=\mathbb C^\beta$, the
viability of a subspace $S\subseteq\mathcal H_A$ for a target
$T\subseteq\mathcal H_A\otimes\mathcal H_B$~\cite[lines
    433--440]{Arad2016RigorousRG} can be stated in $\mathbb C^{\sigma\times\beta}$,
and the cuts inside $\mathcal H_A$ used by the trimming of a viable
subspace~\cite[lines 660--697]{Arad2016RigorousRG} keep $\mathcal H_B$ on the
right.

\begin{theorem}[Cut coordinates with an environment factor]
    \label{thm:cut_environment_coordinates}
    \lean{Entropy.cutEnvEquiv, Entropy.cutEnvVector,
        Entropy.cutEnvVector_leftApply_of_cutOperator_eq_kronecker_one,
        Entropy.cutEnvVector_leftApply_localLift,
        Entropy.cutEnvVector_leftApply_of_cutOperator_eq_one_kronecker,
        Entropy.IsSupportedOn.schmidtRank_cutEnvVector_leftApply_le,
        Entropy.cutEnvVector_rightApply,
        Entropy.schmidtRank_cutEnvVector_rightApply_le,
        Entropy.schmidtRank_cutVector_slice_le}
    \leanok
    \uses{lem:supported_product_decomposition, def:schmidt_rank}
    For a region $D$ of the sites, write a vector of
    $\mathbb C^\sigma\otimes\mathbb C^\beta$ in the bipartite coordinates
    $\mathbb C^{\sigma_D}\otimes(\mathbb C^{\sigma_{D^{\mathrm c}}}\otimes
        \mathbb C^\beta)$; this is an isometry. An operator $K\otimes I$ on the
    sites acts on the left factor, an operator supported outside $D$ or
    acting on $\mathbb C^\beta$ acts on the right factor and does not
    increase the Schmidt rank across $D$, and a slice $w(\cdot,b)$ has
    Schmidt rank across $D$ at most that of $w$.
\end{theorem}

\begin{proof}
    \leanok
    Each statement is a computation of the coefficient matrix in the cut
    coordinates, using the product form of an operator supported in or
    outside $D$; an operator on the right factor multiplies the coefficient
    matrix on the right, which does not increase its rank.
\end{proof}

\begin{theorem}[Viability in coordinates]
    \label{thm:viable_coordinates}
    \lean{Submodule.isCloseTo_map_iff, Submodule.prodExt,
        Submodule.euclideanTensorEquiv, Submodule.map_prodExt_euclideanTensorEquiv,
        Submodule.isCloseTo_prodExt_iff_isViable}
    \leanok
    \uses{def:subspace_close, thm:subspace_close_iff, def:subspace_viable}
    Closeness of subspaces is invariant under linear isometries. Let
    $S\otimes\mathbb C^\beta\subseteq\mathbb C^{\sigma\times\beta}$ be the
    vectors whose slices $w(\cdot,b)$ lie in $S$. The isometry
    $\mathbb C^{\sigma\times\beta}\cong\mathbb C^\sigma\otimes\mathbb C^\beta$
    sending the coordinate vector of $(a,b)$ to $e_a\otimes e_b$ carries
    $S\otimes\mathbb C^\beta$ onto $S\otimes\mathbb C^\beta$ in the tensor
    product, so $S\otimes\mathbb C^\beta$ is $\delta$-close to $T$ exactly
    when $S$ is $\delta$-viable for the image of $T$.
\end{theorem}

\begin{proof}
    \leanok
    Under an isometry $\Phi$, the projection onto $\Phi U$ is $\Phi P_U\Phi^{-1}$,
    so the vector form of closeness transfers. A vector $w$ with slices in $S$
    is $\sum_b w(\cdot,b)\otimes e_b$, and $x\otimes y$ with $x\in S$ has slices
    $y_b\,x\in S$.
\end{proof}

\section{Matrix actions in tensor coordinates}

Let $\mathbb F$ be $\mathbb R$ or $\mathbb C$. For a finite set $X$, write
$\ell^2(X;\mathbb F)$ for its Euclidean space. For finite $X,Y$, let
$\Phi:\ell^2(X;\mathbb F)\otimes\ell^2(Y;\mathbb F)
    \longrightarrow\ell^2(X\times Y;\mathbb F)$ be the unitary linear isometry
which sends the tensor product of the standard basis vectors at $x,y$ to
the standard basis vector at $(x,y)$. These coordinate identities support
finite tensor decompositions as used
in~\cite[Proposition 5.7, lines 1311--1333]{Arad2016RigorousRG}.

\begin{theorem}[The action of a Kronecker matrix]
    \label{thm:euclidean_kronecker_intertwining}
    \lean{Matrix.toEuclideanCLM_kronecker_intertwining}
    \leanok
    If $A$ is a square matrix on $X$ and $B$ is a square matrix on $Y$,
    their Euclidean actions satisfy
    $\Phi\circ(A\otimes B)=(A\otimes_{\mathrm K}B)\circ\Phi$.
    Here the left tensor product is the tensor product of operators and
    $\otimes_{\mathrm K}$ denotes the Kronecker matrix.
\end{theorem}

\begin{proof}\leanok
    In the tensor product of the two standard orthonormal bases, the
    matrix of the tensor product of two linear maps is their Kronecker
    product. Applying the coordinate representation gives the claimed
    equality of continuous linear maps.
\end{proof}

\begin{theorem}[Relabeling a matrix action]
    \label{thm:euclidean_submatrix_intertwining}
    \lean{Matrix.toEuclideanCLM_submatrix_intertwining}
    \leanok
    Let $e:Z\longrightarrow V$ be a bijection between finite sets and let
    $P$ be a square matrix on $V$. Set $P^e_{zz'}=P_{e(z),e(z')}$ and
    $U:\ell^2(Z;\mathbb F)\longrightarrow\ell^2(V;\mathbb F)$ by
    $(Uf)(v)=f(e^{-1}(v))$. Then $U\circ P^e=P\circ U$.
\end{theorem}

\begin{proof}\leanok
    Relabel the summation index in matrix-vector multiplication by $e$.
\end{proof}

\begin{theorem}[A relabeled finite tensor sum]
    \label{thm:euclidean_reindexed_tensor_sum}
    \lean{Matrix.toEuclideanCLM_reindexed_kronecker_sum_intertwining}
    \leanok
    \uses{thm:euclidean_kronecker_intertwining, thm:euclidean_submatrix_intertwining}
    Let $I$ be finite, let $e:Z\longrightarrow X\times Y$ be a bijection,
    and let $A_i,B_i$ be square matrices on $X,Y$, respectively.
    Let $U$ be the relabeling isometry associated to $e$ and put
    $T=\Phi^{-1}\circ U$. Set $P=\sum_{i\in I}A_i\otimes_{\mathrm K}B_i$.
    Then $T\circ P^e=(\sum_{i\in I}A_i\otimes B_i)\circ T$.
    The index set and factor order are unchanged; finite coordinate
    sets and $I$ may be empty.
\end{theorem}

\begin{proof}\leanok
    Sum the Kronecker-action identity over $I$ and combine it with
    relabeling of the matrix action. Apply the injectivity of $\Phi$
    to obtain the equality through $T=\Phi^{-1}\circ U$.
\end{proof}

\begin{theorem}[Exchanging factors in an intertwining decomposition]
    \label{thm:tensor_intertwining_exchange}
    \lean{TensorProduct.intertwining_sum_mapL_commIsometry}
    \leanok
    Let $\mathcal E,\mathcal L,\mathcal R$ be real or complex inner product
    spaces, let $I$ be finite, and let
    $T:\mathcal E\longrightarrow\mathcal L\otimes\mathcal R$ be a unitary
    linear isometry. Suppose $K,A_i,B_i$ are continuous linear operators
    on $\mathcal E,\mathcal L,\mathcal R$, respectively, and
    $TK=(\sum_{i\in I}A_i\otimes B_i)T$.
    Let $S$ exchange the two tensor factors and put $T'=ST$. Then
    \begin{align}
        T'K          & =\Bigl(\sum_{i\in I}B_i\otimes A_i\Bigr)T', \\
        T'K(T')^{-1} & =\sum_{i\in I}B_i\otimes A_i.
    \end{align}
    No finite-dimensionality or completeness assumption is required.
\end{theorem}

\begin{proof}\leanok
    The tensor-factor exchange satisfies
    $S(A_i\otimes B_i)=(B_i\otimes A_i)S$ for every $i$.
    Sum this identity and apply the supplied intertwining relation.
    Evaluation at $(T')^{-1}y$ gives the conjugated identity.
\end{proof}

\begin{theorem}[A conjugated tensor sum with the middle factor first]
    \label{thm:middle_first_conjugated_tensor_sum}
    \lean{Matrix.toEuclideanCLM_middleFirst_conjugation_of_kronecker_sum}
    \leanok
    Let $Z,X,Y,I$ be finite sets and $e:Z\to X\times Y$ a bijection.
    Suppose a matrix $P$ on $Z$ satisfies
    \begin{align}
        P_{zz'}=\sum_{i\in I}A_i(e(z)_X,e(z')_X)B_i(e(z)_Y,e(z')_Y).
    \end{align}
    Let $T:\C^Z\to\C^Y\otimes\C^X$ send the standard vector at $z$ to
    $e_{e(z)_Y}\otimes e_{e(z)_X}$. Then
    \begin{align}
        TPT^{-1}=\sum_{i\in I}B_i\otimes A_i.
    \end{align}
    Empty coordinate and summation sets are allowed.
\end{theorem}
\begin{proof}\leanok
    \uses{thm:euclidean_reindexed_tensor_sum,thm:tensor_intertwining_exchange}
    First apply the reindexed tensor-sum identity with factor order $X,Y$.
    Compose its isometry with exchange of the two tensor factors and then
    with the inverse isometry on the right. No matrix factors commute in
    this argument; only their tensor positions are exchanged.
\end{proof}

\section{Weights of a vector of a tensor product}

\begin{theorem}[Weights along an eigenbasis]
    \label{thm:tensor_weights}
    \lean{Submodule.exists_weights_norm_sq_starProjection_tensorProd}
    \leanok
    \uses{thm:subspace_tensor_projection}
    Let $W$ be a subspace of a finite-dimensional inner product space $M$,
    let $F$ be a finite-dimensional inner product space, and let
    $w\in W\otimes F$. There are an orthonormal basis $f_1,\dots,f_q$ of $W$
    and weights $\mu_k\ge0$ with $\sum_k\mu_k=\lVert w\rVert^2$ such that
    \[
        \lVert P_{W'\otimes F}w\rVert^2=\sum_k\mu_k\lVert P_{W'}f_k\rVert^2
    \]
    for every subspace $W'\subseteq W$. This is the decomposition used in
    the proof of Lemma~l:rs of~\cite[lines 539--541]{Arad2016RigorousRG}.
\end{theorem}

\begin{proof}
    \leanok
    Write $w=\sum_jx_j\otimes e_j$ along an orthonormal basis $e_j$ of $F$;
    then $x_j\in W$ and $\lVert P_{W'\otimes F}w\rVert^2
    =\sum_j\lVert P_{W'}x_j\rVert^2$. Let $f_k$ be an orthonormal eigenbasis
    of $G=\sum_j\lvert x_j\rangle\langle x_j\rvert$ on $W$, with eigenvalues
    $\mu_k\ge0$. Expanding $P_{W'}x_j$ in the basis $f_k$ and using that
    $P_{W'}$ is self-adjoint gives
    $\sum_j\lVert P_{W'}x_j\rVert^2=\sum_k\langle P_{W'}f_k,GP_{W'}f_k\rangle
    =\sum_{k,l}\mu_l\lvert\langle f_l,P_{W'}f_k\rangle\rvert^2
    =\sum_l\mu_l\lVert P_{W'}f_l\rVert^2$. Taking $W'=W$ gives the sum of
    the weights.
\end{proof}

\section{Random sampling of a viable subspace}

Let $M$ and $F$ be finite-dimensional complex inner product spaces, where $F$
stands for $\mathcal H_L\otimes\mathcal H_R$ and $M$ for $\mathcal H_M$ in
the setting of~\cite[lines 532--547]{Arad2016RigorousRG}; viability is read
up to the associator and the exchange of the first two factors.

\begin{theorem}[Net step]
    \label{thm:viable_random_sampling_net}
    \lean{Submodule.isCloseTo_of_net}
    \leanok
    \uses{thm:subspace_close_iff}
    Let $S$ and $T$ be subspaces of a finite-dimensional inner product space
    and let $c\ge0$. If every unit vector $t\in T$ lies within distance
    $\sqrt c$ of some $v\in T$ with $\lVert P_Sv\rVert^2\ge4c$, then
    $\lVert P_St\rVert^2\ge c\lVert t\rVert^2$ for all $t\in T$, that is, $S$
    is $(1-c)$-close to $T$.
\end{theorem}

\begin{proof}
    \leanok
    For a unit vector $u\in T$ choose $v$ as in the hypothesis. Then
    $\lVert P_Su\rVert\ge\lVert P_Sv\rVert-\lVert u-v\rVert
        \ge2\sqrt c-\sqrt c=\sqrt c$. A general nonzero $t\in T$ is
    $\lVert t\rVert$ times a unit vector of $T$.
\end{proof}

\begin{definition}[Random subspace]
    \label{def:viable_random_subspace}
    \lean{Submodule.randomSubspace, Submodule.finrank_randomSubspace}
    \leanok
    Let $W\subseteq M$ have an orthonormal basis indexed by a finite set
    $\iota$, and let $A\subseteq\iota$ have $s$ elements. For a unitary $U$
    on $\mathbb C^\iota$, the subspace $W'_U\subseteq W$ is the image under
    the basis isometry $\mathbb C^\iota\to W$ of $U$ applied to the span of
    the standard basis vectors indexed by $A$; it has dimension $s$. When $U$ is distributed by the
    Haar measure, $W'_U$ is a uniformly random $s$-dimensional subspace of
    $W$.
\end{definition}

\begin{theorem}[One vector]
    \label{thm:viable_random_sampling_one_vector}
    \lean{Submodule.unitaryHaar_norm_sq_starProjection_randomSubspace_le}
    \leanok
    \uses{def:viable_random_subspace, thm:haar_projection_law,
        thm:complex_gaussian_ratio_lower_tail}
    Let $q=\lvert\iota\rvert$. For every unit vector $f\in W$, the Haar
    probability that $\lVert P_{W'_U}f\rVert^2\le s/(2q)$ is at most
    $e^{-s/16}$.
\end{theorem}

\begin{proof}
    \leanok
    The basis isometry carries projections onto subspaces of
    $\mathbb C^\iota$ to projections onto their images, so the event is the
    event that the projection of a unit vector of $\mathbb C^\iota$ onto
    $U$ applied to the coordinate subspace of $A$ has squared norm at most
    $s/(2q)$. Its probability is the probability that
    $\sum_{i\in A}\lvert g_i\rvert^2\le\frac{s}{2q}\sum_i\lvert g_i\rvert^2$
    for independent standard complex Gaussians $g_i$, which is at most
    $e^{-s/16}$.
\end{proof}

\begin{theorem}[Random sampling]
    \label{thm:viable_random_sampling}
    \lean{Submodule.unitaryHaar_not_isViable_randomSubspace_le}
    \leanok
    \uses{thm:viable_random_sampling_net, thm:viable_random_sampling_one_vector,
        thm:tensor_weights, thm:ball_net, def:subspace_viable}
    Let $T\subseteq M\otimes F$ have dimension $r$, let $W\subseteq M$ be
    $\delta$-viable for $T$ with $\delta<1$, and let $W'_U$ be as above with
    $q=\dim W$. Put
    \[
        \delta'=\frac{(1-\delta)s}{8q},\qquad\nu=\sqrt{\delta'},\qquad
        \eta=\Big(1+\frac2\nu\Big)^{2r}qe^{-s/16}.
    \]
    Then the Haar probability that $W'_U$ is not $(1-\delta')$-viable for
    $T$ is at most $\eta$. This is Lemma~l:rs
    of~\cite[lines 532--547]{Arad2016RigorousRG} with the failure
    probability corrected: the printed value
    $(1+4\sqrt{q/((1-\delta)s)})^rqe^{-s/16}$ undercounts the net of the
    unit sphere of $T$, a real space of dimension
    $2r$~\cite{gap:alvv16_random_sampling_net_count}.
\end{theorem}

\begin{proof}
    \leanok
    If $\delta'=0$ every subspace is $1$-viable. Otherwise choose a
    $\nu$-net $C$ of the unit sphere of $T$ with at most $(1+2/\nu)^{2r}$
    points. For $v\in C$ let $w$ be the projection of $v$ onto $W\otimes F$,
    so $\lVert w\rVert^2\ge1-\delta$ and
    $P_{W'_U\otimes F}v=P_{W'_U\otimes F}w$. Write
    $\lVert P_{W'_U\otimes F}w\rVert^2=\sum_k\mu_k\lVert P_{W'_U}f_k\rVert^2$
    along the weights of $w$. Outside the union of the events
    $\lVert P_{W'_U}f_k\rVert^2\le s/(2q)$, over $v\in C$ and the $q$ basis
    vectors $f_k$, every net point satisfies
    $\lVert P_{W'_U\otimes F}v\rVert^2\ge\frac{s}{2q}(1-\delta)=4\delta'$,
    and the net step gives $(1-\delta')$-viability. The union has
    probability at most $\lvert C\rvert qe^{-s/16}\le\eta$.
\end{proof}

\begin{theorem}[Existence of a viable sample]
    \label{thm:viable_random_sampling_exists}
    \lean{Submodule.exists_isViable_randomSubspace}
    \leanok
    \uses{thm:viable_random_sampling}
    In the setting of the random sampling theorem, if $\eta<1$ then some
    unitary $U$ makes $W'_U$ a $(1-\delta')$-viable subspace for $T$ of
    dimension $s$.
\end{theorem}

\begin{proof}
    \leanok
    Otherwise the failure event is the whole unitary group, of probability
    one, which exceeds $\eta$.
\end{proof}

\section{The bootstrapping step}

Let $T,T'\subseteq M\otimes F$ be mutually $0.005$-close, let $D\ge2$, and let
$A_1,\dots,A_{D^2}$ be operators on $M$ such that for every $\delta<1$ and
every $\delta$-viable $S\subseteq M$ for $T'$ the span of $A_1S,\dots,A_{D^2}S$
is $\Delta/(1-\delta)^2$-viable for $T'$, where $D^{12}\Delta\le10^{-5}$. These
are the data of Theorem~t:ais of~\cite[lines 1010--1024]{Arad2016RigorousRG}.
For a dimension $q$ write $s=\lfloor q/(2D^2)\rfloor$ and
\[
    \eta(q)=\Bigl(1+\frac{2}{\sqrt{0.96\,s/(8q)}}\Bigr)^{2\dim T'}q\,e^{-s/16}
\]
for the failure probability of random sampling.

\begin{theorem}[One bootstrapping step]
    \label{thm:viable_bootstrap_step}
    \lean{Submodule.exists_isViable_finrank_le_half}
    \leanok
    \uses{thm:viable_random_sampling, thm:subspace_close_trans}
    If $W\subseteq M$ is $0.015$-viable for $T$, has dimension $q\ge4D^2$, and
    $\eta(q)<1$, then some $0.015$-viable subspace for $T$ has dimension at
    most $q/2$.
\end{theorem}

\begin{proof}
    \uses{thm:finite_image_span_dimension}
    \leanok
    By robustness, $W$ is $0.04$-viable for $T'$. Since $\eta(q)<1$, some
    $s$-dimensional subspace $W'\subseteq W$ is $(1-\delta')$-viable for $T'$
    with $\delta'=0.96\,s/(8q)$. The span $S$ of the $A_iW'$ has dimension at
    most $D^2s\le q/2$ and is $\Delta/\delta'^2$-viable for $T'$. Since
    $s\ge q/(4D^2)$, we have $\delta'\ge0.03/D^2$, and $D^{12}\Delta\le10^{-5}$
    with $D\ge2$ gives $\Delta/\delta'^2\le0.0025$. Robustness again shows
    that $S$ is $2(0.0025+0.005)=0.015$-viable for $T$.
\end{proof}

\begin{theorem}[Bootstrapping]
    \label{thm:viable_bootstrap}
    \lean{Submodule.exists_isViable_finrank_lt_of_bootstrap}
    \leanok
    \uses{thm:viable_bootstrap_step}
    If $q_1\ge4D^2$ and $\eta(q)<1$ for every $q\ge q_1$, then some
    $0.015$-viable subspace for $T$ has dimension less than $q_1$.
    This is the abstract argument used in the proof of
    Proposition~prop:albootstrap-dg of~\cite[lines
        1062--1075]{Arad2016RigorousRG}, once the improving family and the
    sampling bound have been supplied. The spatial scope of the physical
    application and the corrected constants are recorded
    in~\cite{gap:alvv16_bootstrap_constants}.
\end{theorem}

\begin{proof}
    \leanok
    The whole factor $M$ is $0.015$-viable. Starting from it, apply the
    previous theorem as long as the dimension is at least $q_1$; each step
    halves the dimension.
\end{proof}

\begin{theorem}[A quantitative bootstrap bound]
    \label{thm:viable_bootstrap_quantitative_dimension}
    \lean{Submodule.exists_isViable_finrank_lt_rank_mul_pow}
    \leanok
    \uses{def:subspace_close, def:subspace_viable}
    Let $M,F$ be finite-dimensional complex inner-product spaces, and let
    $T,T'\subseteq M\otimes F$ be mutually $0.005$-close. Put
    $r=\dim T'$, and assume $r\ge1$. Let $D\ge2$ be an integer and
    $\Delta\in\R$ satisfy $D^{12}\Delta\le10^{-5}$. Suppose a fixed family
    $A_1,\ldots,A_{D^2}$ of linear operators on $M$ has the following property:
    for every real $\delta<1$ and every $\delta$-viable subspace
    $S\subseteq M$ for $T'$, the space $\operatorname{span}_i A_iS$ is
    $\Delta/(1-\delta)^2$-viable for $T'$.
    Then there is a $0.015$-viable subspace $W\subseteq M$ for $T$ with
    $\dim W<rD^{24}$.

    This is a conditional form of the quantitative argument
    in~\cite[lines 1062--1075]{Arad2016RigorousRG}. Applying it to a spin
    chain requires an improving family for the chosen interval, whose role
    is explained in~\cite{gap:alvv16_bootstrap_constants}.
\end{theorem}
\begin{proof}\leanok
    \uses{thm:viable_bootstrap}
    Let $q\ge rD^{24}$ and $s=\lfloor q/(2D^2)\rfloor$.
    Since $r\ge1$ and $D\ge2$, we have $q\ge4D^2$. The integer quotient
    and the covering-net factor satisfy
    \begin{align}
        s                               & \ge\frac{q}{4D^2}, &
        1+\frac{2}{\sqrt{0.96\,s/(8q)}} & \le21D.
        \label{eq:bootstrap_quantitative_sample}
    \end{align}
    The bound $q\ge rD^{24}$ also gives
    \begin{align}
        \bigl(2r\log(21D)+\log q\bigr)256D^2\le2q.
        \label{eq:bootstrap_quantitative_log_budget}
    \end{align}
    Substituting~(\ref{eq:bootstrap_quantitative_sample})
    and~(\ref{eq:bootstrap_quantitative_log_budget}) into the sampling
    failure probability yields
    \begin{align}
        \eta(q) & \le\exp\bigl(2r\log(21D)+\log q-s/16\bigr)
        \le\exp\bigl(-q/(128D^2)\bigr)<1.
        \label{eq:bootstrap_quantitative_failure}
    \end{align}
    Apply Theorem~\ref{thm:viable_bootstrap} with threshold $rD^{24}$.
\end{proof}

\begin{theorem}[Enlargement at a prescribed accuracy]
    \label{thm:viable_enlargement_inverse_scale}
    \lean{Submodule.viable_enlargement_of_close_improving_target}
    \uses{def:subspace_viable, def:subspace_close}
    Let $\mathcal H,\mathcal M,\mathcal E$ be finite-dimensional complex
    inner-product spaces, and fix an isometry
    $u:\mathcal H\to\mathcal M\otimes\mathcal E$ onto the tensor product.
    Let $G,T\subseteq\mathcal H$, let $1\le m\le D$ and $b\ge0$ be integers,
    and suppose that $G,T$ are mutually $1/(200m)$-close.
    Let $W_0\subseteq\mathcal M$ be $0.015$-viable for $uG$, with $\dim W_0\le b$.
    Suppose that $A_0,\ldots,A_{D^2-1}$ act on $\mathcal M$ and that
    $D^{12}\Delta\le10^{-5}$. Assume that, for every real $\delta<1$ and
    every $\delta$-viable subspace $S$ for $uT$, the space
    $\sum_j A_jS$ is $\Delta/(1-\delta)^2$-viable for $uT$.
    Then $V=\sum_j A_jW_0$ satisfies
    \begin{align}
        \dim V&\le bD^2, & V&\text{ is }\frac1{2m}\text{-viable for }uG.
        \label{eq:viable_inverse_scale}
    \end{align}
    This is a quantitative consequence of the closeness transfers in
    \cite[Lemma~l:close and proof of Theorem~4.9]{Arad2016RigorousRG}.
\end{theorem}
\begin{proof}
    \uses{thm:subspace_close_trans, thm:viable_coordinates, thm:finite_image_span_dimension}
    Since $1/(200m)\le0.005$, the first closeness transfer makes $W_0$
    $0.04$-viable for $uT$. The improving estimate and $m\le D^{12}$ give
    \begin{align*}
        \frac{\Delta}{(1-0.04)^2}&\le\frac1{8m}, &
        2\left(\frac1{8m}+\frac1{200m}\right)&=\frac{13}{50m}\le\frac1{2m}.
    \end{align*}
    For $\Delta\le0$ the first inequality is immediate; for $\Delta>0$
    it follows from $m\Delta\le10^{-5}$.
    Transfer back to $uG$. Finally,
    $\dim\sum_j A_jW_0\le\sum_j\dim A_jW_0\le D^2\dim W_0$.
\end{proof}

\section{Spectral probabilities from viable subspaces}

\begin{theorem}[One spectral distribution for all viable approximations]
    \label{thm:unit_cut_spectral_probability}
    \lean{Submodule.unit_cut_spectral_probability}
    \uses{def:subspace_viable, def:von_neumann_entropy, thm:viable_coordinates}
    Let $\sigma,\beta$ be finite sets, let $\mathcal H$ be a complex
    inner-product space, and fix a linear isometry
    $U:\mathcal H\to\C^\sigma\otimes\C^\beta$ onto the tensor product.
    Let $G\subseteq\mathcal H$ and let $\psi\in G$ be a unit vector.
    Write $U\psi=\sum_{a,b}C_{ab}e_a\otimes e_b$ and $\rho=CC^\dagger$.
    Its eigenvalues $p_0\ge\cdots\ge p_{|\sigma|-1}$ satisfy
    \begin{align*}
        p_a&\ge0, & \sum_a p_a&=1, & S(\rho)&=H(p).
    \end{align*}
    This single distribution has the following property for every real
    $\delta<1$, every subspace $S\subseteq\C^\sigma$, and every integer $R\ge0$:
    if $\dim S\le R$ and $S$ is $\delta$-viable for $UG$, then
    \begin{align}
        \sum_{\substack{0\le a<|\sigma|\\R\le a}}p_a&\le\delta.
        \label{eq:unit_cut_spectral_tail}
    \end{align}
    The distribution is determined before the choices of $\delta,S,R$.
    This is the finite-dimensional spectral step in
    \cite[proof of Theorem~4.9]{Arad2016RigorousRG}.
\end{theorem}
\begin{proof}
    \uses{thm:viable_leading_marginal_mass, thm:entropy_ordered_eigenvalues}
    Positivity of $CC^\dagger$ gives $p_a\ge0$, while
    $\tr(CC^\dagger)=\|\psi\|^2=1$ gives normalization.
    The spectral entropy identity gives $S(\rho)=H(p)$.
    The leading-marginal-mass estimate and monotonicity of the Ky Fan sums give
    $\sum_{a<R}p_a\ge1-\delta$; subtract this inequality from the total mass.
\end{proof}

\begin{theorem}[The actual regional marginal in tensor coordinates]
    \label{thm:actual_region_split_entropy}
    \lean{Submodule.actualRegionSplit_marginal_and_entropy}
    \uses{def:finite_product_regional_entropy, thm:viable_coordinates}
    Let $V$ be a finite set of sites, with local dimensions $n_v\ge0$,
    and let $J\subseteq V$.
    Put $\sigma=\prod_{v\in J}\{0,\ldots,n_v-1\}$ and
    $\beta=\prod_{v\notin J}\{0,\ldots,n_v-1\}$.
    Identify the full configuration space first with
    $\C^\beta\otimes\C^\sigma$ by separating the exterior coordinates,
    and then with $\C^\sigma\otimes\C^\beta$ by exchanging tensor factors.
    For an arbitrary vector $\psi$, let $C$ be its coefficient matrix in
    these coordinates. Then
    \begin{align}
        C_{ab}&=\psi(a\cup b), &
        CC^\dagger&=\tr_{J^c}\ketbra{\psi}{\psi}, &
        S(CC^\dagger)&=S(\rho_J).
        \label{eq:actual_region_split_entropy}
    \end{align}
    Here $a\cup b$ is the full configuration with restrictions $a,b$, and
    $\rho_J=\tr_{J^c}\ketbra{\psi}{\psi}$.
    Neither normalization nor a nonempty region is assumed.
\end{theorem}
\begin{proof}
    \uses{thm:tensor_exchange_product_coordinates, thm:complementary_reduced_ranks}
    Tensor exchange gives $C_{ab}=\psi(a\cup b)$.
    Thus $(CC^\dagger)_{aa'}=\sum_b\psi(a\cup b)\overline{\psi(a'\cup b)}$,
    which is precisely the partial trace. Equality of the matrices gives
    equality of their entropies.
\end{proof}

\section{Error reduction by a spectral projector}

Let $\mathcal H=\mathcal H_L\otimes\mathcal H_M\otimes\mathcal H_R$ and let
$b_j$ be an orthonormal basis of common eigenvectors of a Hamiltonian $H$ and
an operator $K$, with $Hb_j=h_jb_j$ and $Kb_j=k_jb_j$. The low-energy space is
$T=\operatorname{span}\{b_j:h_j\le\eta_0\}$.

\begin{definition}[Spectral projector]
    \label{def:spectral_projector}
    \lean{Submodule.IsSpectralProjector}
    \leanok
    The operator $K$ is a spectral projector for $(H,\eta_0,\eta_1)$ with
    parameter $\Delta$ if $h_j\le\eta_0$ implies $k_j\ge1$ and $h_j\ge\eta_1$
    implies $k_j^2\le\Delta$. This is Definition~def:spectral-AGSP
    of~\cite[lines 570--581]{Arad2016RigorousRG} with closed spectral
    conditions and without positivity of $K$; with the printed open
    conditions the lemma below fails at the eigenvalue
    $\eta_0$~\cite{gap:alvv16_spectral_projector_boundary}.
\end{definition}

\begin{theorem}[Spectral projectors contract the complement]
    \label{thm:spectral_projector_contract}
    \lean{Submodule.IsSpectralProjector.contract}
    \leanok
    \uses{def:spectral_projector}
    If $K$ is a spectral projector and $H$ has no eigenvalue in
    $(\eta_0,\eta_1)$, then $K$ maps $T$ and $T^\perp$ into themselves,
    $\lVert Kt\rVert\ge\lVert t\rVert$ for $t\in T$, and
    $\lVert Kq\rVert^2\le\Delta\lVert q\rVert^2$ for $q\in T^\perp$.
\end{theorem}

\begin{proof}
    \leanok
    In the basis $b_j$, the vectors of $T$ have vanishing coefficients
    outside $\{j:h_j\le\eta_0\}$ and those of $T^\perp$ inside it, and $K$
    multiplies the $j$-th coefficient by $k_j$. Parseval's identity gives
    the two norm bounds.
\end{proof}

\begin{theorem}[Error reduction]
    \label{thm:viable_spectral_projector}
    \lean{Submodule.IsViable.spectralProjector}
    \leanok
    \uses{thm:spectral_projector_contract, thm:subspace_close_contract, thm:subspace_close_mono, def:subspace_viable}
    Let $K=\sum_{i=1}^{D^2}L_i\otimes A_i\otimes R_i$ be a spectral
    projector for $(H,\eta_0,\eta_1)$ with parameter $\Delta\ge0$, let $H$
    have no eigenvalue in $(\eta_0,\eta_1)$, and let $S\subseteq\mathcal H_M$
    be $\delta$-viable for $T$ with $\delta<1$. Then
    $V=\operatorname{span}\{A_iS\}$ has dimension at most $D^2\dim S$ and
    is $\Delta/(1-\delta)^2$-viable for $T$. This is Lemma~l:er
    of~\cite[lines 597--622]{Arad2016RigorousRG}, with the closed spectral
    conditions and the hypothesis $\delta<1$ that its proof uses.
\end{theorem}

\begin{proof}
    \leanok
    The dimension bound holds because $V$ is the sum of the $D^2$ images
    $A_iS$. By the previous theorem, $K$ satisfies the hypotheses of the
    error reduction by a contraction, so
    $K(\mathcal H_L\otimes S\otimes\mathcal H_R)$ is
    $\Delta/(1-\delta)^2$-close to $T$. Each
    $K((l\otimes s)\otimes r)=\sum_i(L_il\otimes A_is)\otimes R_ir$ lies in
    $\mathcal H_L\otimes V\otimes\mathcal H_R$, and closeness passes to the
    larger subspace.
\end{proof}

\section{The Chebyshev filter as a spectral projector}

Let $H$ be a self-adjoint operator on a finite-dimensional inner product space,
with orthonormal eigenbasis $b_j$ and eigenvalues $h_j$.

\begin{theorem}[Polynomials of an operator on its eigenvectors]
    \label{thm:polynomial_on_eigenvector}
    \lean{ContinuousLinearMap.aeval_map_apply_of_apply_eq_smul}
    \leanok
    If $Hv=\mu v$ with $\mu$ real and $p$ is a real polynomial, then
    $p(H)v=p(\mu)v$.
\end{theorem}

\begin{proof}
    \leanok
    Both sides are additive in $p$, and $H^mv=\mu^mv$ by induction on $m$.
\end{proof}

\begin{theorem}[A polynomial filter as a spectral projector]
    \label{thm:polynomial_spectral_projector}
    \lean{Submodule.isSpectralProjector_aeval}
    \leanok
    \uses{thm:polynomial_on_eigenvector, def:spectral_projector}
    If a real polynomial $p$ satisfies $p(t)\ge1$ for $t\le\eta_0$ and
    $p(h_j)^2\le\Delta$ whenever $h_j\ge\eta_1$, then $p(H)$ is a spectral
    projector for $(H,\eta_0,\eta_1)$ with parameter $\Delta$, with eigenvalues
    $p(h_j)$ on the basis $b_j$.
\end{theorem}

\begin{proof}
    \leanok
    The previous theorem gives $p(H)b_j=p(h_j)b_j$, and the two conditions
    are those of the definition.
\end{proof}

\begin{theorem}[The Chebyshev filter]
    \label{thm:chebyshev_spectral_projector}
    \lean{Submodule.exists_isSpectralProjector_chebyshev}
    \leanok
    \uses{thm:polynomial_spectral_projector, thm:chebyshev_interval_polynomial}
    Let all eigenvalues of $H$ be at most $u$, and let $a<b\le u$. For every
    $n$ there is a real polynomial $p$ of degree $n$ with $p(a)=1$ such that
    $p(H)$ is a spectral projector for $(H,a,b)$ with parameter
    $\bigl(2e^{-2n\sqrt{(b-a)/(u-a)}}\bigr)^2$. This is items 1--3 of
    Theorem~thm:Cheby-AGSP of~\cite[lines 789--853]{Arad2016RigorousRG}, in
    the form required by the error reduction; positivity of $p(H)$ is not
    claimed~\cite{gap:alvv16_chebyshev_filter_positivity}.
\end{theorem}

\begin{proof}
    \leanok
    The normalized interval polynomial satisfies $p(t)\ge1$ for $t\le a$ and
    $\lvert p(t)\rvert\le2e^{-2n\sqrt{(b-a)/(u-a)}}$ on $[b,u]$; apply the
    previous theorem.
\end{proof}

\section{Spectral subspaces under isometric coordinates}

\begin{theorem}[Spectral data after an isometry]
    \label{thm:canonical_isometric_spectral_data}
    \lean{Matrix.exists_canonical_isometric_spectral_data}
    \leanok
    Let $B$ be Hermitian on $\C^\Sigma$, with an orthonormal eigenbasis
    $v_i$ and eigenvalues $\lambda_i$. Put
    $G_c=\operatorname{span}\{v_i:\lambda_i\le c\}$.
    Suppose $G_{\eta_0}=G_{\eta_1}$ and $\|B\|\le U$, where
    $\eta_0,\eta_1,U\in\R$. Let $u:\C^\Sigma\to F$ be a unitary linear
    isometry onto a finite-dimensional complex inner-product space.
    The symmetric operator $K=uBu^{-1}$ has an orthonormal
    eigenbasis $b_j$ and eigenvalues $\kappa_j$ satisfying
    \begin{align}
        \kappa_j & \le\eta_0\quad\text{or}\quad\eta_1\le\kappa_j,\notag \\
        \kappa_j & \le U,\notag                                         \\
        uG_c     & =\operatorname{span}\{b_j:\kappa_j\le c\}.
    \end{align}
    The last equality holds for every $c\in\R$. No order between the
    endpoints or nonzero-dimension assumption is needed.
\end{theorem}
\begin{proof}\leanok
    Any orthonormal eigenbasis spans the same closed low spectral subspace:
    expand an eigenvector in another eigenbasis and compare eigenvalues.
    An isometry carries each eigenspace to the eigenspace of the conjugate
    operator. Apply these facts to an orthonormal spectral basis of $K$.
    Equality of the two supplied low spaces excludes intermediate spectral
    values. Finally $\kappa_j\le|\kappa_j|\le\|K\|=\|B\|\le U$.
    These coordinate assertions support the spectral reduction in
    \cite[Theorem 4.2 and Lemma 2.9]{Arad2016RigorousRG}.
\end{proof}

\begin{theorem}[Two spectral values force a nontrivial local dimension]
    \label{thm:two_spectral_values_local_dimension}
    \lean{Matrix.IsHermitian.two_le_local_dimension_of_two_spectral_values}
    \leanok
    Let $H$ be Hermitian on $(\C^q)^{\otimes L}$. If its spectrum contains
    distinct real values $\mu,\nu$, then $q\ge2$.
\end{theorem}
\begin{proof}\leanok
    If $q\le1$, the configuration set has at most one element. Every
    Hermitian matrix on this space has at most one real spectral value,
    contrary to the two given values. This includes the empty configuration
    set and the empty chain.
\end{proof}

\section{Chebyshev improvement from a common polynomial family}

Let $\mathcal H_M$ and $\mathcal H_E$ be finite-dimensional complex inner
product spaces, and let $H$ be self-adjoint on
$\mathcal H=\mathcal H_M\otimes\mathcal H_E$. Choose an orthonormal eigenbasis
$b_j$ with eigenvalues $h_j$. Suppose that $\eta_0<\eta_1\le u$, that
$h_j\le u$, and that each $h_j$ lies in $(-\infty,\eta_0]\cup[\eta_1,\infty)$.
For a nonnegative integer $k$, set
\begin{align}
    G&=\operatorname{span}\{b_j:h_j\le\eta_0\},
    \label{eq:cheb_viable_low_space}\\
    \Delta&=\bigl(2e^{-2k\sqrt{(\eta_1-\eta_0)/(u-\eta_0)}}\bigr)^2.
    \label{eq:cheb_viable_error}
\end{align}
These are the closed spectral conditions used in error reduction; positivity
of the unsquared filter is not required
\cite{gap:alvv16_spectral_projector_boundary,gap:alvv16_chebyshev_filter_positivity}.

\begin{theorem}[Chebyshev improvement from a uniform polynomial reduction]
    \label{thm:chebyshev_viability_polynomial_reduction}
    \lean{Submodule.chebyshev_viability_of_polynomial_reduction}
    \uses{thm:chebyshev_spectral_projector, thm:spectral_projector_contract}
    Fix a nonnegative integer $R$, a map $\Gamma$ on subspaces of
    $\mathcal H_M$, and an operator-valued map $\mathcal F$ on complex
    polynomials such that $\mathcal F(Q)=Q(H)$ for every $Q$.
    Suppose the following reduction holds for every $Q$ of degree at most
    $k$: whenever $S$ is $\delta$-viable for a subspace $T$ of $\mathcal H$,
    $\delta<1$, and $d\ge0$, and $\mathcal F(Q)$ preserves $T$ and $T^\perp$
    and satisfies, for every $t\in T$ and $z\in T^\perp$,
    \begin{align}
        \|\mathcal F(Q)t\|&\ge\|t\|,\\
        \|\mathcal F(Q)z\|^2&\le d\|z\|^2,
    \end{align}
    the subspace $\Gamma(S)$ has dimension at most $R\dim S$ and is
    $d/(1-\delta)^2$-viable for $T$.
    Then, under the spectral assumptions above, for every $\delta$-viable
    subspace $S$ for $G$ with $\delta<1$, the subspace $\Gamma(S)$
    has dimension at most $R\dim S$ and is
    $\Delta/(1-\delta)^2$-viable for $G$.
    Here $\Gamma$ and $\mathcal F$ are fixed before the spectral endpoints
    and before $S$. This conditional composition uses
    Theorem~4.2 of~\cite[lines 789--853]{Arad2016RigorousRG}
    and the error-reduction argument of Lemma~2.9.
\end{theorem}

\begin{proof}
    Choose the normalized real Chebyshev polynomial $p$ of degree $k$ and
    regard it as a complex polynomial. Its evaluation is a spectral
    projector with parameter~(\ref{eq:cheb_viable_error}). The closed spectral
    gap gives invariance, expansion on $G$, and contraction on $G^\perp$.
    Apply the supplied reduction to this polynomial and these bounds.
\end{proof}

\begin{theorem}[A common matrix family gives Chebyshev improvement]
    \label{thm:matrix_family_chebyshev_viability}
    \lean{Matrix.chebyshev_viability_of_conjugated_polynomial_family}
    \uses{thm:viable_tensor_decomposition, thm:chebyshev_viability_polynomial_reduction}
    Let $\Sigma,X,Y$ be finite sets, put $\mathcal H_M=\mathbb C^X$ and
    $\mathcal H_E=\mathbb C^Y$, and fix a unitary linear isometry
    $T:\mathbb C^\Sigma\longrightarrow\mathcal H_M\otimes\mathcal H_E$.
    Fix a matrix $H_0$ on $\mathbb C^\Sigma$ and matrices
    $W_1,\ldots,W_N$ on $\mathbb C^X$. Suppose that for every complex
    polynomial $Q$ of degree at most $k$ there are matrices
    $M_1(Q),\ldots,M_N(Q)$ on $\mathbb C^Y$ such that
    \begin{align}
        TQ(H_0)T^{-1}&=\sum_{r=1}^N W_r\otimes M_r(Q).
        \label{eq:cheb_viable_family}
    \end{align}
    Set $H=TH_0T^{-1}$ and $\Gamma(S)=\sum_{r=1}^N W_rS$.
    If $H$ satisfies the self-adjointness and spectral assumptions stated
    above, then every $\delta$-viable $S$ for
    $G$ of~(\ref{eq:cheb_viable_low_space}), with $\delta<1$, satisfies
    \begin{align}
        \dim\Gamma(S)&\le N\dim S,
        \label{eq:cheb_viable_dimension}\\
        \Gamma(S)&\text{ is }\Delta/(1-\delta)^2\text{-viable for }G.
        \label{eq:cheb_viable_conclusion}
    \end{align}
    The family $(W_r)$ is fixed before the spectral data and before $S$;
    only the complementary matrices $M_r(Q)$ depend on $Q$.
    The supplied polynomial identity is an explicit hypothesis. This is a
    finite-dimensional consequence supporting the polynomial decomposition
    and error-reduction steps of~\cite[Proposition 5.7 and Lemma 2.9]{Arad2016RigorousRG}.
\end{theorem}

\begin{proof}
    Polynomial evaluation commutes with the matrix action on Euclidean
    space and with conjugation by $T$. Thus the filter is $Q(H)$.
    For each $Q$, apply error reduction from its tensor decomposition
    in~(\ref{eq:cheb_viable_family}). The uniform polynomial reduction
    then gives~(\ref{eq:cheb_viable_dimension}) and~(\ref{eq:cheb_viable_conclusion}).
\end{proof}

\begin{theorem}[A numerical bound for the image sum]
    \label{thm:matrix_family_chebyshev_viability_card_bound}
    \lean{Matrix.chebyshev_viability_of_conjugated_polynomial_family_of_card_le}
    \uses{thm:matrix_family_chebyshev_viability}
    Under the hypotheses of the common matrix family theorem, let $C$ be a
    nonnegative integer with $N\le C$. The same fixed subspace
    $\Gamma(S)$ then satisfies $\dim\Gamma(S)\le C\dim S$ and is
    $\Delta/(1-\delta)^2$-viable for $G$, for every $\delta$-viable $S$
    with $\delta<1$.
    The numerical bound $C$ is fixed with the family before the spectral
    endpoints and before $S$.
\end{theorem}

\begin{proof}
    The dimension estimate follows from
    $\dim\Gamma(S)\le N\dim S\le C\dim S$. The viable space and its error
    are unchanged.
\end{proof}

\section{Closeness of low-energy subspaces under a perturbation}

Let $H$ and $H'$ be self-adjoint operators with orthonormal eigenbases
$(b_j)$ and $(b'_j)$ and eigenvalues $a_j$ and $a'_j$. For an energy $c$
write $T(c)$ for the span of the $b_j$ with $a_j\le c$, and $T'(c)$ for the
span of the $b'_j$ with $a'_j\le c$.

\begin{definition}[Low-energy subspace]
    \label{def:spectral_span_le}
    \lean{Submodule.spectralSpanLE}
    \leanok
    The subspace $T(c)$ is the span of the eigenvectors $b_j$ with
    eigenvalue $a_j\le c$.
\end{definition}

\begin{theorem}[Markov's inequality]
    \label{thm:spectral_span_markov}
    \lean{Submodule.mul_norm_starProjection_orthogonal_spectralSpanLE_le}
    \leanok
    \uses{def:spectral_span_le}
    If every eigenvalue $a_j>c$ satisfies $|a_j-\mu|\ge D\ge0$, then the
    projection $P^\perp$ onto the orthogonal complement of $T(c)$ satisfies
    $D\|P^\perp x\|\le\|(H-\mu)x\|$ for every $x$.
\end{theorem}

\begin{proof}
    \leanok
    Expand $x=\sum_jx_jb_j$. Then
    $\|(H-\mu)x\|^2=\sum_j|a_j-\mu|^2|x_j|^2\ge D^2\sum_{a_j>c}|x_j|^2
        =D^2\|P^\perp x\|^2$.
\end{proof}

\begin{theorem}[Eigenvalue comparison from a low-energy residual]
    \label{thm:low_spectral_eigenvalue_comparison}
    \lean{Matrix.IsHermitian.eigenvalues₀_le_add_of_norm_domRestrict_low_le}
    \notready
    \uses{def:spectral_span_le}
    Let $H,F\in\MN{d}$ be Hermitian, with eigenvalues listed in decreasing
    order, and let $G$ be the closed spectral subspace of $F$ at energy $c$.
    Suppose $\|(H-F)|_G\|\le\xi$. For every $0\le i<d$ with
    $\lambda_i(F)\le c$, one has $\lambda_i(H)\le\lambda_i(F)+\xi$.
    This is a variational estimate for use with the restricted perturbation
    bounds in~\cite[lines 955--968]{Arad2016RigorousRG}.
\end{theorem}
\begin{proof}
    \notready
    The first $i+1$ eigenvectors of $H$ and the last $d-i$ eigenvectors of
    $F$ span subspaces with a common nonzero vector $x$. The latter subspace
    is contained in $G$. Spectral expansion and the residual estimate give
    \begin{align}
        \lambda_i(H)\|x\|^2
         & \le \operatorname{Re}\langle Hx,x\rangle
        \le \operatorname{Re}\langle Fx,x\rangle+\xi\|x\|^2
        \le (\lambda_i(F)+\xi)\|x\|^2.
        \notag
    \end{align}
    Divide by $\|x\|^2>0$.
\end{proof}

\begin{theorem}[One window]
    \label{thm:spectral_span_window}
    \lean{Submodule.mul_norm_starProjection_orthogonal_spectralSpanLE_le_of_window}
    \leanok
    \uses{thm:spectral_span_markov}
    Suppose $\|(H-H')\varphi\|\le\xi\|\varphi\|$ for every $\varphi\in T(c)$,
    and let $h,D\ge0$ with $\mu+D\le c'$. If $\psi\in T(c)$ is supported on
    eigenvectors $b_j$ with $|a_j-\mu|\le h$, then the projection $P'^\perp$
    onto the orthogonal complement of $T'(c')$ satisfies
    $D\|P'^\perp\psi\|\le(\xi+h)\|\psi\|$.
\end{theorem}

\begin{proof}
    \leanok
    Every eigenvalue $a'_k>c'$ satisfies $|a'_k-\mu|\ge D$. By the previous
    theorem,
    $D\|P'^\perp\psi\|\le\|(H'-\mu)\psi\|\le\|(H'-H)\psi\|
        +\|(H-\mu)\psi\|\le(\xi+h)\|\psi\|$.
\end{proof}

\begin{theorem}[Closeness of the low-energy subspaces]
    \label{thm:spectral_span_close_of_perturbation}
    \lean{Submodule.isCloseTo_spectralSpanLE_of_perturbation}
    \leanok
    \uses{thm:spectral_span_window, def:subspace_close}
    Suppose $\|(H-H')\varphi\|\le\xi\|\varphi\|$ for every $\varphi\in T(c)$,
    every eigenvalue $a_j\le c$ is at least $m$, $m<c\le c'$ and $h>0$. Let
    $N=\lceil(c-m)/(2h)\rceil$. If $N(\xi+h)^2\le\delta(c'-c+h)^2$, then
    $T'(c')$ is $\delta$-close to $T(c)$. This is the closeness step in the
    proof of the second item of Lemma lem:hard-trunc
    of~\cite[lines 953--969]{Arad2016RigorousRG}, with the number of windows
    corrected~\cite{gap:alvv16_hard_truncation_windows}.
\end{theorem}

\begin{proof}
    \leanok
    Cover $[m,c]$ by the $N$ closed windows of width $2h$ centred at
    $c-h-2hw$ for $w<N$. Split $\varphi\in T(c)$ as
    $\varphi=\sum_{w<N}\psi_w$, where $\psi_w$ keeps the coefficients of the
    $b_j$ with $\lceil(c-a_j)/(2h)\rceil-1=w$, counting $a_j=c$ in the top
    window. Each $\psi_w$ lies in $T(c)$ and is supported within $h$ of its
    centre, and every centre plus $c'-c+h$ is at most $c'$. By the previous
    theorem and the Cauchy--Schwarz inequality,
    \[
        (c'-c+h)\|P'^\perp\varphi\|\le(\xi+h)\sum_w\|\psi_w\|
        \le(\xi+h)\sqrt{N}\,\|\varphi\|,
    \]
    so $\|P'^\perp\varphi\|^2\le\delta\|\varphi\|^2$ and
    $\|P'\varphi\|^2\ge(1-\delta)\|\varphi\|^2$.
\end{proof}

\begin{theorem}[A low-energy cluster from restricted residual estimates]
    \label{thm:low_spectral_cluster_restricted_residuals}
    \lean{Submodule.low_spectral_cluster_of_restricted_residuals}
    \leanok
    \uses{def:spectral_span_le, def:subspace_close}
    Let $H,H'$ act on a finite-dimensional complex inner-product space,
    with complete orthonormal eigenbases and real eigenvalues $a_j,a'_j$.
    Suppose $\gamma>0$, $0\le\xi\le\gamma/10^6$, and every $a_j$ equals
    $e$ or is at least $e+\gamma$. Assume the two restricted operator
    estimates
    \begin{align}
        \bigl\|(H-H')|_{T(e+\gamma/40)}\bigr\|    & \le\xi,
        \label{eq:low_cluster_original_residual}            \\
        \bigl\|(H-H')|_{T'(e+19\gamma/20)}\bigr\| & \le\xi.
        \label{eq:low_cluster_perturbed_residual}
    \end{align}
    Put $G=T(e)$, $U=T'(e+\gamma/20)$ and $V=T'(e+19\gamma/20)$.
    Then $U=V$ and $\dim U=\dim G$. Each of $U$ and $G$ is
    $1/200$-close to the other. For every
    $e+\gamma/20\le\tau\le e+19\gamma/20$, the space $T'(\tau)$ equals
    $U$ and retains both approximation bounds. Every perturbed eigenvalue
    satisfies
    \begin{align}
        a'_j\le e+\gamma/10
        \quad\text{or}\quad e+9\gamma/10\le a'_j.
        \label{eq:low_cluster_eigenvalue_gap}
    \end{align}
    The cutoffs are absolute energies relative to $e$; no comparison
    with the perturbed minimum energy is asserted. This is a
    finite-dimensional consequence of the estimates
    in~\cite[lines 955--968 and 1010--1024]{Arad2016RigorousRG}, using the
    corrected window count~\cite{gap:alvv16_hard_truncation_windows}.
\end{theorem}

\begin{proof}
    \leanok
    \uses{thm:spectral_span_close_of_perturbation,
        thm:subspace_close_nested_equal, thm:subspace_close_dimension,
        thm:spectral_span_quadratic_form}
    The original spectral dichotomy gives $H\ge eI$ and
    $T(e+\gamma/40)=T(e+39\gamma/40)=G$. By
    (\ref{eq:low_cluster_perturbed_residual}), a unit eigenvector $x$ of
    $H'$ with eigenvalue $a'_j\le e+19\gamma/20$ satisfies
    $e\le\operatorname{Re}\langle x,Hx\rangle\le a'_j+\xi$.
    Thus these perturbed eigenvalues are at least $e-\xi$.

    Apply the low-energy perturbation estimate with
    $c=e+\gamma/40$, $c'=e+\gamma/20$, $m=e$ and then with the two
    operators interchanged and
    $c=e+19\gamma/20$, $c'=e+39\gamma/40$, $m=e-\xi$.
    In each application take $h=\gamma/10^6$. The source width is at
    most $\gamma$, the cutoff separation is $\gamma/40$, and the number
    of windows is at most $500001$. Consequently
    \begin{align}
        N(\xi+h)^2
         & \le 500001(2\gamma/10^6)^2
        \le \frac1{200}(\gamma/40)^2
        \le \frac1{200}(c'-c+h)^2.
        \label{eq:low_cluster_window_bound}
    \end{align}
    The two applications show that $U$ is $1/200$-close to $G$ and $G$
    is $1/200$-close to $V$. Since $U\subseteq V$, the nested-subspace
    theorem gives $U=V$, equality of dimensions with $G$, and both
    approximation bounds for every intermediate cutoff. If
    $e+\gamma/20<a'_j\le e+19\gamma/20$, its eigenvector belongs to
    $V=U$, whereas its own coordinate must vanish in $U$. This
    contradiction proves~(\ref{eq:low_cluster_eigenvalue_gap}).
\end{proof}


\begin{theorem}[Energy on a low spectral subspace]
    \label{thm:spectral_span_quadratic_form}
    \lean{OrthonormalBasis.re_inner_apply_eq_sum_of_eigen,
        Submodule.re_inner_apply_le_of_mem_spectralSpanLE}
    \leanok
    \uses{def:spectral_span_le}
    Let $(b_j)$ be a finite orthonormal basis of a complex inner-product
    space, and let $H$ be a linear operator satisfying $Hb_j=a_jb_j$ with
    $a_j\in\R$. For every $x$,
    \begin{align}
        \operatorname{Re}\langle x,Hx\rangle
         & =\sum_j a_j|\langle b_j,x\rangle|^2.
        \label{eq:spectral_span_energy_coordinates}
    \end{align}
    In particular, if $x\in T(c)$, then
    $\operatorname{Re}\langle x,Hx\rangle\le c\|x\|^2$.
\end{theorem}
\begin{proof}
    \leanok
    Express $x$ and $Hx$ in the orthonormal basis to obtain
    (\ref{eq:spectral_span_energy_coordinates}). For $x\in T(c)$, all
    coordinates with $a_j>c$ vanish. Bound the remaining coefficients by
    $c$ and apply Parseval's identity.
\end{proof}

\begin{theorem}[The low spectral projection commutes with the operator]
    \label{thm:spectral_span_projection_commutes}
    \lean{Submodule.commute_starProjection_spectralSpanLE}
    \leanok
    \uses{def:spectral_span_le}
    Under the same eigenbasis assumptions, the orthogonal projection
    onto $T(c)$ commutes with $H$, for every real $c$.
\end{theorem}
\begin{proof}
    \leanok
    In that basis the complementary projection retains precisely the
    coordinates with $a_j>c$, while $H$ multiplies each coordinate by $a_j$.
    These two operations commute, and so do $H$ and the low projection.
    These are spectral consequences used
    in~\cite[lines 955--968 and 1010--1024]{Arad2016RigorousRG}.
\end{proof}

\begin{theorem}[Relative low-energy spaces above the actual ground energies]
    \label{thm:relative_low_energy_spaces_close}
    \lean{Matrix.IsHermitian.groundEnergy_compare_and_isCloseTo_spectralSpanLE_of_norm_domRestrict_le}
    \notready
    \uses{def:spectral_span_le,def:subspace_close}
    Let $H,F$ be Hermitian matrices on a nonzero finite-dimensional space,
    with $F\le H$. Let $E$ and $E'$ be their respective least eigenvalues,
    and write $T_H(c)$ and $T_F(c)$ for their closed spectral subspaces
    at the absolute energy $c$.
    Suppose $\eta,\xi,\delta>0$ and
    \begin{align}
        \|(H-F)|_{T_H(E+\eta)}\| & \le\xi, &
        \|(H-F)|_{T_F(E+\eta)}\| & \le\xi.
        \label{eq:relative_low_energy_residuals}
    \end{align}
    Put $N=\lceil\eta/(2\xi)\rceil$ and
    $\eta'=\eta+\xi+2\xi\sqrt{N/\delta}$.
    Then $E'\le E\le E'+\xi$,
    $T_F(E'+\eta')$ is $\delta$-close to $T_H(E+\eta)$, and
    $T_H(E+\eta')$ is $\delta$-close to $T_F(E'+\eta)$.
    Both relative cutoffs use the actual respective ground energy.
    This is the finite-dimensional relative-energy consequence of
    \cite[lines 955--968]{Arad2016RigorousRG}, with the window count and
    ground-energy shift corrected in
    \cite{gap:alvv16_hard_truncation_windows}.
\end{theorem}
\begin{proof}
    \notready
    \uses{thm:low_spectral_eigenvalue_comparison,
        thm:spectral_span_close_of_perturbation}
    Weyl monotonicity gives $E'\le E$. Since the lowest eigenvalue of
    $F$ lies below $E+\eta$, the restricted eigenvalue comparison gives
    $E\le E'+\xi$.
    Apply the window estimate with width $h=\xi$ to the spaces of $H$
    between $E$ and $E+\eta$. For the enlarged target cutoff, the
    denominator satisfies
    \begin{align}
        (E'+\eta')-(E+\eta)+\xi
         & \ge\xi+2\xi\sqrt{N/\delta}
        \ge2\xi\sqrt{N/\delta}.
        \notag
    \end{align}
    Its squared value times $\delta$ is at least $N(2\xi)^2$, so the
    finite-window budget holds. The target cutoff is at least the source
    cutoff by the same ground-energy comparison.
    For the reverse direction, $E'\le E$ gives
    $T_F(E'+\eta)\subseteq T_F(E+\eta)$, so the second residual estimate
    remains valid on the smaller space. Its spectral width is again
    $\eta$, and the reverse target denominator is larger still.
    Apply the same window argument with $H,F$ interchanged.
\end{proof}

\section{The detectability lemma and its converse}

Let $h_0,\dots,h_{m-1}$ be orthogonal projections on a complete inner
product space, write $Q_j=1-h_j$, and for a vector $\psi$ put
$\psi_k=Q_{k-1}\cdots Q_0\psi$. The detectability lemma bounds the energy
of $\varphi=\psi_m$ with respect to $H=\sum_jh_j$ by the mass that the
product $Q_{m-1}\cdots Q_0$ removes from $\psi$. It is Lemma~lem:dl of
\cite[lines 1508--1519]{Arad2016RigorousRG}, proved in
\cite{Anshu2016Simple}; its converse is Lemma~lem:dl-converse of
\cite[lines 1528--1537]{Arad2016RigorousRG}.

\begin{definition}[Product of complementary projections]
    \label{def:complement_product}
    \lean{ContinuousLinearMap.complementProd}
    \leanok
    The product $Q_{k-1}\cdots Q_0$, in which $Q_0=1-h_0$ acts first.
\end{definition}

\begin{theorem}[Projection identities]
    \label{thm:star_projection_identities}
    \lean{ContinuousLinearMap.IsStarProjection.apply_apply_self,
        ContinuousLinearMap.IsStarProjection.inner_apply_left,
        ContinuousLinearMap.IsStarProjection.inner_apply_self,
        ContinuousLinearMap.IsStarProjection.re_inner_apply_self,
        ContinuousLinearMap.IsStarProjection.norm_sq_eq,
        ContinuousLinearMap.IsStarProjection.norm_apply_le}
    \leanok
    For an orthogonal projection $p$ and vectors $x,y$:
    $p(px)=px$, $\langle px,y\rangle=\langle x,py\rangle$,
    $\langle x,px\rangle=\|px\|^2$, $\|x\|^2=\|px\|^2+\|x-px\|^2$ and
    $\|px\|\le\|x\|$.
\end{theorem}

\begin{proof}
    \leanok
    The first two hold because $p$ is idempotent and self-adjoint, and they
    give $\langle x,px\rangle=\langle px,px\rangle$ and
    $\langle px,x-px\rangle=0$.
\end{proof}

\begin{theorem}[Removed mass]
    \label{thm:complement_product_removed_mass}
    \lean{ContinuousLinearMap.complementProd_succ_apply,
        ContinuousLinearMap.norm_sq_complementProd_succ,
        ContinuousLinearMap.re_inner_sum_apply_eq,
        ContinuousLinearMap.sum_norm_sq_apply_complementProd}
    \leanok
    \uses{def:complement_product, thm:star_projection_identities}
    For every $k$, $\psi_{k+1}=\psi_k-h_k\psi_k$ and
    $\|\psi_{k+1}\|^2=\|\psi_k\|^2-\|h_k\psi_k\|^2$. Hence
    $\sum_{k<n}\|h_k\psi_k\|^2=\|\psi\|^2-\|\psi_n\|^2$. Moreover
    $\langle y,(\sum_{i\in s}h_i)y\rangle=\sum_{i\in s}\|h_iy\|^2$ for every
    finite set $s$ and vector $y$.
\end{theorem}

\begin{proof}
    \leanok
    Apply Pythagoras to $\psi_k=h_k\psi_k+Q_k\psi_k$ and sum over $k$.
\end{proof}

\begin{theorem}[One term of the energy]
    \label{thm:detectability_one_term}
    \lean{ContinuousLinearMap.norm_apply_complementProd_le}
    \leanok
    \uses{thm:complement_product_removed_mass}
    If $i<n$, then $\|h_i\psi_n\|\le\sum_k\|h_k\psi_k\|$, the sum over the
    indices $i<k<n$ for which $h_k$ does not commute with $h_i$.
\end{theorem}

\begin{proof}
    \leanok
    By induction on $n$. For $n=i+1$, $h_i\psi_{i+1}=h_iQ_i\psi_i=0$. If
    $h_n$ commutes with $h_i$, then $h_i\psi_{n+1}=Q_nh_i\psi_n$, whose norm
    is at most $\|h_i\psi_n\|$. Otherwise
    $h_i\psi_{n+1}=h_i\psi_n-h_ih_n\psi_n$, whose norm is at most
    $\|h_i\psi_n\|+\|h_n\psi_n\|$.
\end{proof}

\begin{theorem}[Detectability lemma]
    \label{thm:detectability_lemma}
    \lean{ContinuousLinearMap.re_inner_sum_complementProd_le}
    \leanok
    \uses{thm:detectability_one_term}
    Suppose each $h_i$ fails to commute with at most $g$ of the other
    $h_j$, $j<m$. Then
    $\langle\varphi|H|\varphi\rangle\le g^2(\|\psi\|^2-\|\varphi\|^2)$.
\end{theorem}

\begin{proof}
    \leanok
    Let $A_i$ be the set of $k$ with $i<k<m$ for which $h_k$ does not
    commute with $h_i$, and $a_k=\|h_k\psi_k\|$. By the previous theorem
    and the Cauchy--Schwarz inequality,
    $\|h_i\varphi\|^2\le|A_i|\sum_{k\in A_i}a_k^2\le g\sum_{k\in A_i}a_k^2$.
    Each $k$ lies in $A_i$ for at most $g$ indices $i$, so
    $\langle\varphi|H|\varphi\rangle=\sum_i\|h_i\varphi\|^2
        \le g^2\sum_{k<m}a_k^2=g^2(\|\psi\|^2-\|\varphi\|^2)$.
\end{proof}

\begin{theorem}[Detectability lemma, normalized form]
    \label{thm:detectability_lemma_normalized}
    \lean{ContinuousLinearMap.norm_sq_complementProd_le}
    \leanok
    \uses{thm:detectability_lemma}
    In the setting of the previous theorem, let $\|\psi\|=1$ and
    $\varepsilon_\varphi=\langle\varphi|H|\varphi\rangle/\|\varphi\|^2$.
    Then $\|\varphi\|^2\le1/(\varepsilon_\varphi/g^2+1)$, where a quotient
    by zero is read as zero.
\end{theorem}

\begin{proof}
    \leanok
    If $\varphi\ne0$ and $g>0$, then
    $\|\varphi\|^2(\varepsilon_\varphi/g^2+1)
        =\langle\varphi|H|\varphi\rangle/g^2+\|\varphi\|^2\le1$. If
    $\varphi=0$ there is nothing to prove, and if $g=0$ the right side is
    $1\ge\|\varphi\|^2$.
\end{proof}

\begin{theorem}[Commuting projections]
    \label{thm:commuting_complement_product}
    \lean{ContinuousLinearMap.commute_complementProd,
        ContinuousLinearMap.isStarProjection_complementProd,
        ContinuousLinearMap.norm_sq_sub_norm_sq_complementProd_le}
    \leanok
    \uses{thm:complement_product_removed_mass}
    An operator that commutes with $h_0,\dots,h_{k-1}$ commutes with
    $Q_{k-1}\cdots Q_0$. If the $h_j$ commute pairwise, then
    $Q_{n-1}\cdots Q_0$ is an orthogonal projection and
    $\|\psi\|^2-\|\psi_n\|^2\le\langle\psi|\sum_{k<n}h_k|\psi\rangle$.
\end{theorem}

\begin{proof}
    \leanok
    A product of commuting orthogonal projections is an orthogonal
    projection. For the inequality, $h_k\psi_k=Q_{k-1}\cdots Q_0h_k\psi$
    has norm at most $\|h_k\psi\|$, and
    $\|h_k\psi\|^2=\langle\psi|h_k|\psi\rangle$.
\end{proof}

\begin{theorem}[Product of two projections]
    \label{thm:two_projection_product}
    \lean{ContinuousLinearMap.IsStarProjection.one_sub_three_mul_le_norm_sq_apply_apply}
    \leanok
    \uses{thm:star_projection_identities}
    For orthogonal projections $P,P'$ and a unit vector $x$,
    $\|PP'x\|^2\ge1-3(\|x-Px\|^2+\|x-P'x\|^2)$.
\end{theorem}

\begin{proof}
    \leanok
    Put $u=x-Px$ and $w=x-P'x$. Then
    $\langle x,PP'x\rangle=\langle x-u,x-w\rangle$, whose real part is
    $1-\|u\|^2-\|w\|^2+\operatorname{Re}\langle u,w\rangle
        \ge1-\tfrac32(\|u\|^2+\|w\|^2)$. This real part is at most
    $\|PP'x\|$, and squaring gives the claim when
    $1-\tfrac32(\|u\|^2+\|w\|^2)\ge0$; otherwise the right side is negative.
\end{proof}

\begin{theorem}[Converse of the detectability lemma]
    \label{thm:detectability_lemma_converse}
    \lean{ContinuousLinearMap.one_sub_three_mul_le_norm_sq_complementProd_comp}
    \leanok
    \uses{thm:commuting_complement_product, thm:two_projection_product}
    Let $(h_k)$ and $(h'_k)$ be two families of pairwise commuting
    orthogonal projections, $Q_k=1-h_k$, $Q'_k=1-h'_k$,
    $P=Q_{n-1}\cdots Q_0$ and $P'=Q'_{n'-1}\cdots Q'_0$, and $H=\sum_{k<n}h_k+\sum_{k<n'}h'_k$. Then
    $\|PP'\psi\|^2\ge1-3\langle\psi|H|\psi\rangle$ for every unit vector
    $\psi$.
\end{theorem}

\begin{proof}
    \leanok
    By the previous two theorems with $\|\psi-P\psi\|^2=1-\|P\psi\|^2$.
\end{proof}

For a nearest-neighbour chain whose terms are projectors, the even and the
odd terms form two such families, and $PP'$ is the operator $\mathrm{DL}(H)$ of
\cite[Definition def:DLO]{Arad2016RigorousRG}. The converse above holds for
every unit vector, not only for eigenvectors of $H$, and its bound
$1-3\langle\psi|H|\psi\rangle$ implies the printed bound
$1-4\langle\psi|H|\psi\rangle$.

\begin{theorem}[Order and kernels under range projection]
    \label{thm:support_projection_order_kernel}
    \lean{Matrix.IsHermitian.supportProj_eq_cfc}
    \lean{Matrix.IsHermitian.isStarProjection_supportProj}
    \lean{Matrix.PosSemidef.le_supportProj_of_norm_le_one}
    \lean{Matrix.PosSemidef.sum_supportProj_mulVec_eq_zero_iff}
    \leanok
    For a finite Hermitian matrix $A$, its range projection is
    $P_A=\mathbf1_{\mathbb R\setminus\{0\}}(A)$ and is an orthogonal projection.
    If $A\ge0$ and $\|A\|\le1$, then $A\le P_A$.
    For every finite family $h_i\ge0$,
    $\ker\sum_i h_i=\ker\sum_i P_{h_i}$.
    This is the range-projection reduction used in
    \cite[Section 6.1]{Arad2016RigorousRG}.
\end{theorem}
\begin{proof}\leanok
    In an orthonormal eigenbasis, $P_A$ has diagonal entries zero or one.
    On the spectrum of a positive contraction, $x\le\mathbf1_{x\ne0}$.
    Functional calculus gives the order comparison. A positive sum annihilates
    a vector precisely when every summand does. Each positive matrix and its
    range projection have the same kernel, proving the last identity.
\end{proof}

\section{Trimming at a threshold on the Schmidt coefficients}

Arad, Landau, Vazirani and Vidick trim a viable subspace cut by cut, keeping
at each cut the left Schmidt vectors whose Schmidt coefficient is at least
$\xi$~\cite[lines 642--656]{Arad2016RigorousRG}. For a bipartite vector
$\psi\in\mathbb C^m\otimes\mathbb C^n$ with coefficient matrix $C$ and left
marginal $\rho=CC^\dagger$, the squared Schmidt coefficients are the
eigenvalues of $\rho$, so the kept vectors span the range of the spectral
projection $P$ of $\rho$ onto $[\xi^2,\infty)$. This section proves the facts
about one such projection used in the proof of their trimming
lemma~\cite[lines 660--697]{Arad2016RigorousRG}.

\begin{theorem}[Pairing with small Schmidt coefficients]
    \label{thm:schmidt_threshold_pairing}
    \lean{Matrix.norm_inner_le_of_posSemidef_smul_one_sub_leftMarginal}
    \leanok
    \uses{def:schmidt_rank}
    Let $c\ge0$, let $\varphi$ have left marginal at most $c$ in the Loewner
    order, and let $\psi$ have Schmidt rank at most $k$. Then
    $|\langle\varphi,\psi\rangle|\le\sqrt{ck}\,\|\psi\|$.
\end{theorem}

\begin{proof}
    \leanok
    Let $R$ be the support projection of the left marginal of $\psi$; it has
    rank at most $k$ and $(R\otimes I)\psi=\psi$. Hence
    $\langle\varphi,\psi\rangle=\langle(R\otimes I)\varphi,\psi\rangle$, and
    $\|(R\otimes I)\varphi\|^2=\tr(R\rho_\varphi R)\le c\tr R\le ck$.
\end{proof}

\begin{theorem}[The discarded part]
    \label{thm:schmidt_threshold_discarded}
    \lean{Matrix.PosSemidef.posSemidef_smul_one_sub_one_sub_spectralProjectionGE_mul,
        Matrix.posSemidef_smul_one_sub_leftMarginal_leftApply_one_sub,
        Matrix.norm_inner_leftApply_one_sub_spectralProjectionGE_le}
    \leanok
    \uses{def:closed_threshold_projection, thm:closed_threshold_projection,
        thm:schmidt_threshold_pairing}
    Let $b\ge0$ and let $P$ be the spectral projection of the left marginal
    $\rho$ of $\psi$ onto $[b,\infty)$. Then $(I-P)\rho\le b$, so the
    discarded part $((I-P)\otimes I)\psi$ has left marginal at most $b$, and
    for every $\varphi$ of Schmidt rank at most $k$,
    \[
        |\langle((I-P)\otimes I)\psi,\varphi\rangle|\le\sqrt{bk}\,\|\varphi\|.
    \]
    With $b=\xi^2$ this is the bound $\xi\sqrt k$ of
    \cite[lines 666--680]{Arad2016RigorousRG}.
\end{theorem}

\begin{proof}
    \leanok
    On an eigenvalue $t\ge b$ of $\rho$ the function $(1-[t\ge b])\,t$
    vanishes, and on $t<b$ it equals $t<b$. The left marginal of the
    discarded part is $(I-P)\rho(I-P)=(I-P)\rho$, since $P$ commutes with
    $\rho$. The pairing bound is
    Theorem~\ref{thm:schmidt_threshold_pairing}.
\end{proof}

\begin{theorem}[The kept part]
    \label{thm:schmidt_threshold_kept}
    \lean{Matrix.mul_rank_spectralProjectionGE_leftMarginal_le,
        Matrix.exists_mul_eq_mul_transpose_of_commute,
        Matrix.exists_leftApply_spectralProjectionGE_eq_rightApply}
    \leanok
    \uses{def:closed_threshold_projection, thm:closed_threshold_projection,
        thm:trimming_second_factor}
    With $P$ as above, $b\operatorname{rank}P\le\|\psi\|^2$. Moreover, every
    orthogonal projection $P$ commuting with $\rho$ satisfies $PC=CQ^{\mathsf T}$
    for an orthogonal projection $Q$ of rank at most $\operatorname{rank}P$,
    that is $(P\otimes I)\psi=(I\otimes Q)\psi$. In particular applying $P$
    does not increase the Schmidt rank across a cut inside the second factor
    \cite[lines 690--692]{Arad2016RigorousRG}.
\end{theorem}

\begin{proof}
    \leanok
    Each eigenvalue of $\rho$ counted by $P$ is at least $b$, and their sum is
    at most $\tr\rho=\|\psi\|^2$. The second statement is the argument of
    Theorem~\ref{thm:trimming_second_factor}, which uses only that $P$ is an
    orthogonal projection commuting with $\rho$.
\end{proof}

\section{Trimming a viable subspace}

Arad, Landau, Vazirani and Vidick reduce the Schmidt rank of a viable
subspace by trimming it collectively, cut by cut~\cite[lines
    642--697]{Arad2016RigorousRG}. Here $\mathcal H_A=\mathbb C^\sigma$ is the
space of configurations of a finite set of sites, the cuts are nested regions
$D_0\subseteq D_1\subseteq\cdots$, the cut $j$ putting the sites of $D_j$ on
the left, and the environment $\mathcal H_B=\mathbb C^\beta$ lies on the right
of every cut. For a chain
$\mathcal H_A=\mathcal H_A^1\otimes\cdots\otimes\mathcal H_A^\ell$ the region
$D_j$ consists of the first $j+1$ factors, and there are $\ell-1$ cuts.

\begin{definition}[Trimming]
    \label{def:viable_trimming}
    \lean{Entropy.stackVector, Entropy.trimProj, Entropy.trimStep,
        Entropy.trimSubspace}
    \leanok
    \uses{def:closed_threshold_projection, thm:cut_environment_coordinates}
    Let $u_1,\dots,u_s$ be an orthonormal basis of $S\subseteq\mathcal H_A$ and
    $\psi=\sum_iu_i\otimes|i\rangle$. Put $Q_0=I$ and, for $j\ge0$, let
    $P^{j+1}$ be the spectral projection onto $[\xi^2,\infty)$ of the left
    marginal, across the cut $D_j$, of $(Q_j\otimes I)\psi$, and
    $Q_{j+1}=Q_j(P^{j+1}\otimes I)$. The trimmed subspace after $m$ cuts is
    $\operatorname{Trim}_\xi(S)=\operatorname{span}\{Q_mu_i\}$. This is
    Definition~def:trimming of~\cite[lines 642--656]{Arad2016RigorousRG}.
\end{definition}

\begin{theorem}[Nesting of the trimming projections]
    \label{thm:viable_trimming_nested}
    \lean{Matrix.mul_spectralProjectionGE_leftMarginal_of_leftApply_eq,
        Entropy.trimProj_props, Entropy.trimProj_mul_trimProj_of_le,
        Entropy.trimDiff, Entropy.trimDiff_eq,
        Entropy.sum_norm_leftApply_trimDiff_sq_le}
    \leanok
    \uses{def:viable_trimming}
    For $\xi>0$ each $Q_j$ is an orthogonal projection supported in $D_j$,
    $Q_{j+1}=P^{j+1}\otimes I$, and $Q_kQ_j=Q_jQ_k=Q_j$ for $k\le j$. The
    differences $Q_j-Q_{j+1}$ are mutually orthogonal projections, so
    $\sum_{j<m}\|((Q_j-Q_{j+1})\otimes I)t\|^2\le\|t\|^2$.
\end{theorem}

\begin{proof}
    \leanok
    If $(M\otimes I)\varphi=\varphi$ then $M$ fixes the left marginal $\rho$ of
    $\varphi$, and the spectral projection onto $[b,\infty)$ with $b>0$ is
    $\rho\,g(\rho)$ with $g(t)=t^{-1}$ there, so $M$ fixes it as well. Apply
    this to $M$ the lift of $Q_j$ across the cut $D_j$, which fixes
    $(Q_j\otimes I)\psi$. The orthogonality is then a telescoping Pythagoras
    identity.
\end{proof}

\begin{theorem}[Trimming]
    \label{thm:viable_trimming}
    \lean{Matrix.inner_leftApply_left, Matrix.inner_rightApply_left,
        Matrix.norm_rightApply_sq_le, Entropy.eq_rightApply_stackVector,
        Entropy.sum_norm_inner_sq_eq, Entropy.norm_inner_leftApply_trimDiff_le,
        Entropy.leftApply_trimProj_mem_prodExt,
        Entropy.isCloseTo_prodExt_trimSubspace,
        Entropy.schmidtRank_cutVector_trimmed_le}
    \leanok
    \uses{thm:viable_trimming_nested, thm:schmidt_threshold_discarded,
        thm:schmidt_threshold_kept, thm:viable_coordinates, thm:subspace_close_iff}
    Let $S\otimes\mathbb C^\beta$ be $\delta$-close to $T$, let every vector of
    $T$ have Schmidt rank at most $b$ across each of the cuts
    $D_0,\dots,D_{m-1}$, and let $\xi>0$. Then
    $\operatorname{Trim}_\xi(S)\otimes\mathbb C^\beta$ is
    $(\delta+2\xi\sqrt{mb})$-close to $T$, and each $Q_mu_i$ has Schmidt rank
    at most $s\xi^{-2}$ across each cut. This is
    Lemma~lem:trimming of~\cite[lines 660--697]{Arad2016RigorousRG}, with the
    bound on $\delta'$ that its argument gives; the printed
    $\delta+\sqrt{\ell bs}\,\xi$ does not follow from that
    argument~\cite{gap:alvv16_trimming_constant}.
\end{theorem}

\begin{proof}
    \leanok
    For $t\in T$ let $w=\sum_iu_i\otimes c_i$ be its projection onto
    $S\otimes\mathbb C^\beta$. The part of $w$ removed at the cut $j$ pairs with
    $t$ as the discarded part of $\psi$ at that cut pairs with
    $W=(I\otimes C^\dagger)w_j$, where $w_j$ is the piece of $t$ removed at the
    cut and $C=(c_i(b))$. The vector $W$ has Schmidt rank at most $b$ and norm
    at most $\|w\|\,\|w_j\|$, so the pairing is at most
    $\xi\sqrt b\,\|w\|\,\|w_j\|$. Summing over the cuts with
    $\sum_j\|w_j\|\le\sqrt m\,\|t\|$ gives
    $\|P\,t\|\ge\|w\|-\xi\sqrt{mb}\,\|t\|$ for the projection $P$ onto the
    trimmed extension, and squaring gives the bound. For the ranks, $Q_m$ lies
    below $P^{k+1}\otimes I$, whose rank is at most $\|\psi\|^2\xi^{-2}=s\xi^{-2}$.
\end{proof}

